
\Data\ID{1003027701}
\Updated{2020-07-15 03:07:12}
\Level{A}
\SubArea{Definite integral}
\LANGen{
\Question{
How much bigger is \( \int\limits_0^{\pi}(x+\sin x)\,\mathrm{d}x \) than \( \int\limits_1^{\pi}x\,\mathrm{d}x \)?}
{\( \frac52 \)}{\( 0.5 \)}{\( \pi-1 \)}{\( \frac32 \)}{}{}}
\LANGpl{
\Question{
O ile jest większa  \( \int\limits_0^{\pi}(x+\sin x)\,\mathrm{d}x \) od \( \int\limits_1^{\pi}x\,\mathrm{d}x \)?}
{\( \frac52 \)}{\( 0{,}5 \)}{\( \pi-1 \)}{\( \frac32 \)}{}{}}
\LANGcs{
\Question{
O kolik je větší \( \int\limits_0^{\pi}(x+\sin x)\,\mathrm{d}x \) než \( \int\limits_1^{\pi}x\,\mathrm{d}x \)?}
{o \( \frac52 \)}{o \( 0{,}5 \)}{o \( \pi-1 \)}{o \( \frac32 \)}{}{}}
\LANGsk{
\Question{
O koľko je väčší \( \int\limits_0^{\pi}(x+\sin x)\,\mathrm{d}x \) než \( \int\limits_1^{\pi}x\,\mathrm{d}x \)?}
{o \( \frac52 \)}{o \( 0{,}5 \)}{o \( \pi-1 \)}{o \( \frac32 \)}{}{}}
\LANGes{
\Question{
¿Cuántas veces más grande es la integral \( \int\limits_0^{\pi}(x+\sin x)\,\mathrm{d}x \) con respecto a \( \int\limits_1^{\pi}x\,\mathrm{d}x \)?}
{\( \frac52 \)}{\( 0.5 \)}{\( \pi-1 \)}{\( \frac32 \)}{}{}}
\End

\Data\ID{1003027702}
\Updated{2020-07-15 03:07:12}
\Level{A}
\SubArea{Definite integral}
\LANGen{
\Question{
How much smaller is \( \int\limits_{-\frac{\pi}2}^{\frac{\pi}2}\sin x\,\mathrm{d}x \) than \( \int\limits_{0}^{\pi}(\cos x + 1)\,\mathrm{d}x \)?}
{\( \pi \)}{\( 1-\pi \)}{\( \pi-1 \)}{\( 2\pi \)}{}{}}
\LANGpl{
\Question{
O ile mniejsza jest  \( \int\limits_{-\frac{\pi}2}^{\frac{\pi}2}\sin x\,\mathrm{d}x \) od \( \int\limits_{0}^{\pi}(\cos x + 1)\,\mathrm{d}x \)?}
{\( \pi \)}{\( 1-\pi \)}{\( \pi-1 \)}{\( 2\pi \)}{}{}}
\LANGcs{
\Question{
O kolik je menší  \( \int\limits_{-\frac{\pi}2}^{\frac{\pi}2}\sin x\,\mathrm{d}x \) než \( \int\limits_{0}^{\pi}(\cos x + 1)\,\mathrm{d}x \)?}
{o \( \pi \)}{o \( 1-\pi \)}{o \( \pi-1 \)}{o \( 2\pi \)}{}{}}
\LANGsk{
\Question{
O koľko je menší  \( \int\limits_{-\frac{\pi}2}^{\frac{\pi}2}\sin x\,\mathrm{d}x \) než \( \int\limits_{0}^{\pi}(\cos x + 1)\,\mathrm{d}x \)?}
{o \( \pi \)}{o \( 1-\pi \)}{o \( \pi-1 \)}{o \( 2\pi \)}{}{}}
\LANGes{
\Question{
¿Cuánto más pequeña es \( \int\limits_{-\frac{\pi}2}^{\frac{\pi}2}\sin x\,\mathrm{d}x \) con relación a \( \int\limits_{0}^{\pi}(\cos x + 1)\,\mathrm{d}x \)?}
{\( \pi \)}{\( 1-\pi \)}{\( \pi-1 \)}{\( 2\pi \)}{}{}}
\End

\Data\ID{1003027703}
\Updated{2020-07-15 03:07:12}
\Level{A}
\SubArea{Definite integral}
\LANGen{
\Question{
Compare two definite integrals \( I_1=\int\limits_0^{\frac{\pi}2}\cos x\,\mathrm{d}x \) and \( I_2=\int\limits_0^{2\pi}2\cos x\,\mathrm{d}x \).}
{\( I_1 \) is bigger than \( I_2 \) by \( 1 \).}{\( I_1 \) is smaller than \( I_2 \) by \( 1 \).}{\( I_1 \) is equal to \( I_2 \).}{\( I_1 \) is smaller than \( I_2 \) by \( 2 \).}{}{}}
\LANGpl{
\Question{
Porównaj dwie całki oznaczone \( I_1=\int\limits_0^{\frac{\pi}2}\cos x\,\mathrm{d}x \) i  \( I_2=\int\limits_0^{2\pi}2\cos x\,\mathrm{d}x \).}
{\( I_1 \) jest większa od \( I_2 \) o \( 1 \).}{\( I_1 \) jest mniejsza od  \( I_2 \) o \( 1 \).}{\( I_1 \) jest równa \( I_2 \).}{\( I_1 \) jest mniejsza od \( I_2 \) o \( 2 \).}{}{}}
\LANGcs{
\Question{
Porovnejte dva určité integrály \( I_1=\int\limits_0^{\frac{\pi}2}\cos x\,\mathrm{d}x \) a \( I_2=\int\limits_0^{2\pi}2\cos x\,\mathrm{d}x \).}
{\( I_1 \) je větší než \( I_2 \) o \( 1 \).}{\( I_1 \) je menší než \( I_2 \) o \( 1 \).}{\( I_1 \) je roven \( I_2 \).}{\( I_1 \) je menší než \( I_2 \) o \( 2 \).}{}{}}
\LANGsk{
\Question{
Porovnajte dva určité integrály \( I_1=\int\limits_0^{\frac{\pi}2}\cos x\,\mathrm{d}x \) a \( I_2=\int\limits_0^{2\pi}2\cos x\,\mathrm{d}x \).}
{\( I_1 \) je väčší než \( I_2 \) o \( 1 \).}{\( I_1 \) je menší než \( I_2 \) o \( 1 \).}{\( I_1 \) je rovný \( I_2 \).}{\( I_1 \) je menší než \( I_2 \) o \( 2 \).}{}{}}
\LANGes{
\Question{
Compara dos integrales definidas \( I_1=\int\limits_0^{\frac{\pi}2}\cos x\,\mathrm{d}x \) y \( I_2=\int\limits_0^{2\pi}2\cos x\,\mathrm{d}x \).}
{\( I_1 \) es más grande que \( I_2 \) de \( 1 \).}{\( I_1 \) es más pequeña que \( I_2 \) de \( 1 \).}{\( I_1 \) es igual a \( I_2 \).}{\( I_1 \) es más pequeña que \( I_2 \) de \( 2 \).}{}{}}
\End

\Data\ID{1003027704}
\Updated{2020-07-15 04:07:14}
\Level{A}
\SubArea{Definite integral}
\LANGen{
\Question{
Compare two definite integrals \( I_1=\int\limits_{-1}^1x^8\,\mathrm{d}x \) and \( I_2=\int\limits_{-1}^1x^2\,\mathrm{d}x \).}
{\( I_2 \) is bigger than \( I_1 \) by \( \frac49 \).}{\( I_1 \) is bigger than \( I_2 \) by \( \frac49 \).}{\( I_2 \) is bigger than \( I_1 \) by \( \frac29 \).}{\( I_1 \) is bigger than \( I_2 \) by \( \frac29 \).}{}{}}
\LANGpl{
\Question{
Porównaj dwie całki oznaczone \( I_1=\int\limits_{-1}^1x^8\,\mathrm{d}x \) i \( I_2=\int\limits_{-1}^1x^2\,\mathrm{d}x \).}
{\( I_2 \) jest większa niż \( I_1 \) o \( \frac49 \).}{\( I_1 \) jest większa niż  \( I_2 \) o \( \frac49 \).}{\( I_2 \) jest większa niż \( I_1 \) o \( \frac29 \).}{\( I_1 \) jest większa niż \( I_2 \) o \( \frac29 \).}{}{}}
\LANGcs{
\Question{
Porovnejte dva určité integrály \( I_1=\int\limits_{-1}^1x^8\,\mathrm{d}x \) a \( I_2=\int\limits_{-1}^1x^2\,\mathrm{d}x \).}
{\( I_2 \) je větší než \( I_1 \) o \( \frac49 \).}{\( I_1 \) je větší než \( I_2 \) o \( \frac49 \).}{\( I_2 \) je větší než \( I_1 \) o \( \frac29 \).}{\( I_1 \) je větší než \( I_2 \) o \( \frac29 \).}{}{}}
\LANGsk{
\Question{
Porovnajte dva určité integrály \( I_1=\int\limits_{-1}^1x^8\,\mathrm{d}x \) a \( I_2=\int\limits_{-1}^1x^2\,\mathrm{d}x \).}
{\( I_2 \) je väčší než \( I_1 \) o \( \frac49 \).}{\( I_1 \) je väčší než \( I_2 \) o \( \frac49 \).}{\( I_2 \) je väčší než \( I_1 \) o \( \frac29 \).}{\( I_1 \) je väčší než \( I_2 \) o \( \frac29 \).}{}{}}
\LANGes{
\Question{
Compara dos integrales definidas \( I_1=\int\limits_{-1}^1x^8\,\mathrm{d}x \) y \( I_2=\int\limits_{-1}^1x^2\,\mathrm{d}x \).}
{\( I_2 \) es más grande que \( I_1 \) de \( \frac49 \).}{\( I_1 \) es más grande que \( I_2 \) de \( \frac49 \).}{\( I_2 \) es más grande que \( I_1 \) de \( \frac29 \).}{\( I_1 \) es más grande que \( I_2 \) de \( \frac29 \).}{}{}}
\End

\Data\ID{1003027705}
\Updated{2020-07-15 03:07:12}
\Level{A}
\SubArea{Definite integral}
\LANGen{
\Question{
How much smaller is \( \int\limits_0^{1} x^{15}\,\mathrm{d}x \) than  \( \int\limits_0^{1} x^{5}\,\mathrm{d}x \)?}
{\( \frac5{48} \)}{\( \frac{11}{48} \)}{\( \frac5{32} \)}{\( \frac{11}{32} \)}{}{}}
\LANGpl{
\Question{
O ile mniejsza jest \( \int\limits_0^{1} x^{15}\,\mathrm{d}x \) od  \( \int\limits_0^{1} x^{5}\,\mathrm{d}x \)?}
{\( \frac5{48} \)}{\( \frac{11}{48} \)}{\( \frac5{32} \)}{\( \frac{11}{32} \)}{}{}}
\LANGcs{
\Question{
O kolik je menší \( \int\limits_0^{1} x^{15}\,\mathrm{d}x \) než  \( \int\limits_0^{1} x^{5}\,\mathrm{d}x \)?}
{o \( \frac5{48} \)}{o \( \frac{11}{48} \)}{o \( \frac5{32} \)}{o \( \frac{11}{32} \)}{}{}}
\LANGsk{
\Question{
O koľko je menší \( \int\limits_0^{1} x^{15}\,\mathrm{d}x \) než  \( \int\limits_0^{1} x^{5}\,\mathrm{d}x \)?}
{o \( \frac5{48} \)}{o \( \frac{11}{48} \)}{o \( \frac5{32} \)}{o \( \frac{11}{32} \)}{}{}}
\LANGes{
\Question{
¿Cuánto más pequeña es \( \int\limits_0^{1} x^{15}\,\mathrm{d}x \) con respecto a \( \int\limits_0^{1} x^{5}\,\mathrm{d}x \)?}
{\( \frac5{48} \)}{\( \frac{11}{48} \)}{\( \frac5{32} \)}{\( \frac{11}{32} \)}{}{}}
\End

\Data\ID{1003027706}
\Updated{2020-07-15 03:07:12}
\Level{A}
\SubArea{Definite integral}
\LANGen{
\Question{
How much smaller is \( \int\limits_{\frac{5\pi}6}^{\pi}3\sin x\,\mathrm{d}x \) than \( \int\limits_0^{\frac{5\pi}6}3\sin x\,\mathrm{d}x \)?}
{\( 3\sqrt3 \)}{\( 3\sqrt2 \)}{6}{\( \frac{\pi}6 \)}{}{}}
\LANGpl{
\Question{
O ile mniejsza jest wartość \( \int\limits_{\frac{5\pi}6}^{\pi}3\sin x\,\mathrm{d}x \) od \( \int\limits_0^{\frac{5\pi}6}3\sin x\,\mathrm{d}x \)?}
{\( 3\sqrt3 \)}{\( 3\sqrt2 \)}{6}{\( \frac{\pi}6 \)}{}{}}
\LANGcs{
\Question{
O kolik je menší \( \int\limits_{\frac{5\pi}6}^{\pi}3\sin x\,\mathrm{d}x \) než \( \int\limits_0^{\frac{5\pi}6}3\sin x\,\mathrm{d}x \)?}
{o \( 3\sqrt3 \)}{o \( 3\sqrt2 \)}{o 6}{o \( \frac{\pi}6 \)}{}{}}
\LANGsk{
\Question{
O koľko je menší \( \int\limits_{\frac{5\pi}6}^{\pi}3\sin x\,\mathrm{d}x \) než \( \int\limits_0^{\frac{5\pi}6}3\sin x\,\mathrm{d}x \)?}
{o \( 3\sqrt3 \)}{o \( 3\sqrt2 \)}{o 6}{o \( \frac{\pi}6 \)}{}{}}
\LANGes{
\Question{
¿Cuánto más pequeña es \( \int\limits_{\frac{5\pi}6}^{\pi}3\sin x\,\mathrm{d}x \) con respecto a \( \int\limits_0^{\frac{5\pi}6}3\sin x\,\mathrm{d}x \)?}
{\( 3\sqrt3 \)}{\( 3\sqrt2 \)}{6}{\( \frac{\pi}6 \)}{}{}}
\End

\Data\ID{1003027707}
\Updated{2020-07-15 03:07:12}
\Level{A}
\SubArea{Definite integral}
\LANGen{
\Question{
Which of the given values of \( a \) makes the statement true?
\[ \int\limits_0^51.6x\,\mathrm{d}x-a = \int\limits_0^50.6x\,\mathrm{d}x \]}
{\( a=12.5 \)}{\( a=1 \)}{\( a=13.5 \)}{\( a=8 \)}{}{}}
\LANGpl{
\Question{
Dla których wartości  \( a \) zachodzi równanie?
\[ \int\limits_0^51{,}6x\,\mathrm{d}x-a = \int\limits_0^50{,}6x\,\mathrm{d}x \]}
{\( a=12{,}5 \)}{\( a=1 \)}{\( a=13{,}5 \)}{\( a=8 \)}{}{}}
\LANGcs{
\Question{
Pro kterou z daných hodnot \( a\) platí tato rovnost?
\[ \int\limits_0^51{,}6x\,\mathrm{d}x-a = \int\limits_0^50{,}6x\,\mathrm{d}x \]}
{\( a=12{,}5 \)}{\( a=1 \)}{\( a=13{,}5 \)}{\( a=8 \)}{}{}}
\LANGsk{
\Question{
Pre ktorú z daných hodnôt \( a\) platí táto rovnosť?
\[ \int\limits_0^51{,}6x\,\mathrm{d}x-a = \int\limits_0^50{,}6x\,\mathrm{d}x \]}
{\( a=12{,}5 \)}{\( a=1 \)}{\( a=13{,}5 \)}{\( a=8 \)}{}{}}
\LANGes{
\Question{
¿Cuál de los valores dados de \( a \) hace el enunciado correcto?
\[ \int\limits_0^51.6x\,\mathrm{d}x-a = \int\limits_0^50.6x\,\mathrm{d}x \]}
{\( a=12.5 \)}{\( a=1 \)}{\( a=13.5 \)}{\( a=8 \)}{}{}}
\End

\Data\ID{1003027708}
\Updated{2020-07-15 03:07:12}
\Level{A}
\SubArea{Definite integral}
\LANGen{
\Question{
How many times bigger is \( \int\limits_{5}^{10}(1.2x-6)\,\mathrm{d}x \) than \( \int\limits_{0}^{5}0.4x\,\mathrm{d}x \)?}
{\( 3 \) times}{\( 2 \) times}{\( 4 \) times}{\( 7 \) times}{}{}}
\LANGpl{
\Question{
Ile razy wartość \( \int\limits_{5}^{10}(1{,}2x-6)\,\mathrm{d}x \) jest większa od \( \int\limits_{0}^{5}0{,}4x\,\mathrm{d}x \)?}
{\( 3 \) razy}{\( 2 \) razy}{\( 4 \) razy}{\( 7 \) razy}{}{}}
\LANGcs{
\Question{
Kolikrát je větší \( \int\limits_{5}^{10}(1{,}2x-6)\,\mathrm{d}x \) než \( \int\limits_{0}^{5}0{,}4x\,\mathrm{d}x \)?}
{třikrát}{dvakrát}{čtyřikrát}{sedmkrát}{}{}}
\LANGsk{
\Question{
Koľkokrát je väčší \( \int\limits_{5}^{10}(1{,}2x-6)\,\mathrm{d}x \) než \( \int\limits_{0}^{5}0{,}4x\,\mathrm{d}x \)?}
{trikrát}{dvakrát}{štyrikrát}{sedemkrát}{}{}}
\LANGes{
\Question{
¿Cuánto más grande es \( \int\limits_{5}^{10}(1.2x-6)\,\mathrm{d}x \) con respecto a \( \int\limits_{0}^{5}0.4x\,\mathrm{d}x \)?}
{\( 3 \) veces}{\( 2 \) veces}{\( 4 \) veces}{\( 7 \) veces}{}{}}
\End

\Data\ID{1003027709}
\Updated{2020-07-15 03:07:12}
\Level{A}
\SubArea{Definite integral}
\LANGen{
\Question{
How many times bigger is \( \int\limits_0^7\left(-\frac47x+4\right)\,\mathrm{d}x \) than \( \int\limits_{-1}^0\left(-\frac43x+\frac43\right)\,\mathrm{d}x \)?}
{\( 7 \) times}{\( 6 \) times}{\( 14 \) times}{\( 2 \) times}{}{}}
\LANGpl{
\Question{
Ile razy wartość \( \int\limits_0^7\left(-\frac47x+4\right)\,\mathrm{d}x \) jest większa od  \( \int\limits_{-1}^0\left(-\frac43x+\frac43\right)\,\mathrm{d}x \)?}
{\( 7 \) razy}{\( 6 \) razy}{\( 14 \) razy}{\( 2 \) razy}{}{}}
\LANGcs{
\Question{
Kolikrát je větší \( \int\limits_0^7\left(-\frac47x+4\right)\,\mathrm{d}x \) než \( \int\limits_{-1}^0\left(-\frac43x+\frac43\right)\,\mathrm{d}x \)?}
{sedmkrát}{šestkrát}{čtrnáctkrát}{dvakrát}{}{}}
\LANGsk{
\Question{
Koľkokrát je väčší \( \int\limits_0^7\left(-\frac47x+4\right)\,\mathrm{d}x \) než \( \int\limits_{-1}^0\left(-\frac43x+\frac43\right)\,\mathrm{d}x \)?}
{sedemkrát}{šesťkrát}{štrnásťkrát}{dvakrát}{}{}}
\LANGes{
\Question{
¿Cuánto más grande es \( \int\limits_0^7\left(-\frac47x+4\right)\,\mathrm{d}x \) con respecto a \( \int\limits_{-1}^0\left(-\frac43x+\frac43\right)\,\mathrm{d}x \)?}
{\( 7 \) veces}{\( 6 \) veces}{\( 14 \) veces}{\( 2 \) veces}{}{}}
\End

\Data\ID{1003027710}
\Updated{2020-07-15 03:07:12}
\Level{A}
\SubArea{Definite integral}
\LANGen{
\Question{
How many times bigger is \( \int\limits_0^4 (3x+3)\,\mathrm{d}x \) than \( \int\limits_2^3 (3x-6)\,\mathrm{d}x \)?}
{\( 24 \) times}{\( 6 \) times}{\( 9 \) times}{\( 54 \) times}{}{}}
\LANGpl{
\Question{
Ile razy wartość  \( \int\limits_0^4 (3x+3)\,\mathrm{d}x \) jest większa od  \( \int\limits_2^3 (3x-6)\,\mathrm{d}x \)?}
{\( 24 \) razy}{\( 6 \) razy}{\( 9 \) razy}{\( 54 \) razy}{}{}}
\LANGcs{
\Question{
Kolikrát je větší \( \int\limits_0^4 (3x+3)\,\mathrm{d}x \) než \( \int\limits_2^3 (3x-6)\,\mathrm{d}x \)?}
{dvacetčtyřikrát}{šestkrát}{devětkrát}{padesátčtyřikrát}{}{}}
\LANGsk{
\Question{
Koľkokrát je väčší \( \int\limits_0^4 (3x+3)\,\mathrm{d}x \) než \( \int\limits_2^3 (3x-6)\,\mathrm{d}x \)?}
{dvadsaťštyrikrát}{šesťkrát}{deväťkrát}{päťdesiatštyrikrát}{}{}}
\LANGes{
\Question{
¿Cuántas veces más grande es \( \int\limits_0^4 (3x+3)\,\mathrm{d}x \) con respecto a \( \int\limits_2^3 (3x-6)\,\mathrm{d}x \)?}
{\( 24 \) veces}{\( 6 \) veces}{\( 9 \) veces}{\( 54 \) veces}{}{}}
\End

\Data\ID{1003027711}
\Updated{2020-07-15 03:07:12}
\Level{A}
\SubArea{Definite integral}
\LANGen{
\Question{
Compare two definite integrals \( I_1 = \int\limits_0^5 0.6x\,\mathrm{d}x \) and \( I_2 = \int\limits_0^5 1.8x\,\mathrm{d}x \).}
{\( I_2 \) is \( 3 \) times as big as \( I_1 \).}{\( I_1 \) is \( 3 \) times as big as \( I_2 \).}{\( I_2 \) is \( 1.2 \) multiple of \( I_1 \).}{\( I_2 \) is \( 30 \) times as big as \( I_ 1 \).}{}{}}
\LANGpl{
\Question{
Porównaj dwie całki oznaczone  \( I_1 = \int\limits_0^5 0{,}6x\,\mathrm{d}x \) i \( I_2 = \int\limits_0^5 1{,}8x\,\mathrm{d}x \).}
{\( I_2 \) jest \( 3 \) razy większa od \( I_1 \).}{\( I_1 \) jest  \( 3 \) razy większa od \( I_2 \).}{\( I_2 \) jest \( 1{,}2 \) wielokrotnością \( I_1 \).}{\( I_2 \) jest \( 30 \) razy większa od \( I_ 1 \).}{}{}}
\LANGcs{
\Question{
Porovnejte dva určité integrály \( I_1 = \int\limits_0^5 0{,}6x\,\mathrm{d}x \) a \( I_2 = \int\limits_0^5 1{,}8x\,\mathrm{d}x \).}
{\( I_2 \) je třikrát větší než \( I_1 \).}{\( I_1 \) je třikrát větší než \( I_2 \).}{\( I_2 \) je \( 1{,}2 \) násobkem \( I_1 \).}{\( I_2 \) je třicetkrát větší než \( I_ 1 \).}{}{}}
\LANGsk{
\Question{
Porovnajte dva určité integrály \( I_1 = \int\limits_0^5 0{,}6x\,\mathrm{d}x \) a \( I_2 = \int\limits_0^5 1{,}8x\,\mathrm{d}x \).}
{\( I_2 \) je trikrát väčší než \( I_1 \).}{\( I_1 \) je trikrát väčší než \( I_2 \).}{\( I_2 \) je \( 1{,}2 \) násobkom \( I_1 \).}{\( I_2 \) je tridsaťkrát väčší než \( I_ 1 \).}{}{}}
\LANGes{
\Question{
Compara dos integrales definidas \( I_1 = \int\limits_0^5 0.6x\,\mathrm{d}x \) y \( I_2 = \int\limits_0^5 1.8x\,\mathrm{d}x \).}
{\( I_2 \) es \( 3 \) veces más grande que \( I_1 \).}{\( I_1 \) es \( 3 \) veces más grande que \( I_2 \).}{\( I_2 \) es \( 1.2 \) múltiplo de \( I_1 \).}{\( I_2 \) es \( 30 \) veces más grande que \( I_ 1 \).}{}{}}
\End

\Data\ID{1003027712}
\Updated{2020-07-15 03:07:12}
\Level{A}
\SubArea{Definite integral}
\LANGen{
\Question{
Compare two definite integrals \(  I_1 = \int\limits_1^2 \frac1x\,\mathrm{d}x \) and \( I_2 = \int\limits_2^4 \frac1x\,\mathrm{d}x \).}
{\( I_1 \) is equal to \( I_2 \).}{\( I_2 \) is twice as big as \( I_1 \).}{\( I_1 \) is twice as big as  \( I_2 \).}{\( I_1 \) is \( 4 \) times as big as \( I_2 \).}{}{}}
\LANGpl{
\Question{
Porównaj dwie całki oznaczone  \(  I_1 = \int\limits_1^2 \frac1x\,\mathrm{d}x \) i \( I_2 = \int\limits_2^4 \frac1x\,\mathrm{d}x \).}
{\( I_1 \) jest równa \( I_2 \).}{\( I_2 \) jest dwa razy większa od \( I_1 \).}{\( I_1 \) jest dwa razy większa od \( I_2 \).}{\( I_1 \) jest \( 4 \) razy większa od  \( I_2 \).}{}{}}
\LANGcs{
\Question{
Porovnejte dva určité integrály \(  I_1 = \int\limits_1^2 \frac1x\,\mathrm{d}x \) a \( I_2 = \int\limits_2^4 \frac1x\,\mathrm{d}x \).}
{\( I_1 \) je roven \( I_2 \).}{\( I_2 \) je dvojnásobkem \( I_1 \).}{\( I_1 \) je dvojnásobkem \( I_2 \).}{\( I_1 \) je čtyřnásobkem \( I_2 \).}{}{}}
\LANGsk{
\Question{
Porovnajte dva určité integrály \(  I_1 = \int\limits_1^2 \frac1x\,\mathrm{d}x \) a \( I_2 = \int\limits_2^4 \frac1x\,\mathrm{d}x \).}
{\( I_1 \) je rovný \( I_2 \).}{\( I_2 \) je dvojnásobkom \( I_1 \).}{\( I_1 \) je dvojnásobkom \( I_2 \).}{\( I_1 \) je štvornásobkom \( I_2 \).}{}{}}
\LANGes{
\Question{
Compara dos integrales definidas \(  I_1 = \int\limits_1^2 \frac1x\,\mathrm{d}x \) y \( I_2 = \int\limits_2^4 \frac1x\,\mathrm{d}x \).}
{\( I_1 \) es igual a \( I_2 \).}{\( I_2 \) es dos veces más grande que \( I_1 \).}{\( I_1 \) es dos veces más grande que \( I_2 \).}{\( I_1 \) es \( 4 \) veces más grande que \( I_2 \).}{}{}}
\End

\Data\ID{1003027713}
\Updated{2020-07-15 03:07:12}
\Level{A}
\SubArea{Definite integral}
\LANGen{
\Question{
Which of the given values of \( a \) makes the statement true?
\[ a\cdot\int\limits_0^{\frac{\pi}2}\cos x\,\mathrm{d}x=\int\limits_{2\pi}^{3\pi}\sin x\,\mathrm{d}x \]}
{\( a=2 \)}{\( a=\pi \)}{\( a=4 \)}{\( a=1 \)}{}{}}
\LANGpl{
\Question{
Dla jakiej wartości \( a \) zachodzi równość?
\[ a\cdot\int\limits_0^{\frac{\pi}2}\cos x\,\mathrm{d}x=\int\limits_{2\pi}^{3\pi}\sin x\,\mathrm{d}x \]}
{\( a=2 \)}{\( a=\pi \)}{\( a=4 \)}{\( a=1 \)}{}{}}
\LANGcs{
\Question{
Pro kterou z daných hodnot \( a \) platí rovnost?
\[ a\cdot\int\limits_0^{\frac{\pi}2}\cos x\,\mathrm{d}x=\int\limits_{2\pi}^{3\pi}\sin x\,\mathrm{d}x \]}
{\( a=2 \)}{\( a=\pi \)}{\( a=4 \)}{\( a=1 \)}{}{}}
\LANGsk{
\Question{
Pre ktorú z daných hodnôt \( a \) platí rovnosť?
\[ a\cdot\int\limits_0^{\frac{\pi}2}\cos x\,\mathrm{d}x=\int\limits_{2\pi}^{3\pi}\sin x\,\mathrm{d}x \]}
{\( a=2 \)}{\( a=\pi \)}{\( a=4 \)}{\( a=1 \)}{}{}}
\LANGes{
\Question{
¿Cuál de los valores dados de \( a \) hace el enunciado correcto?
\[ a\cdot\int\limits_0^{\frac{\pi}2}\cos x\,\mathrm{d}x=\int\limits_{2\pi}^{3\pi}\sin x\,\mathrm{d}x \]}
{\( a=2 \)}{\( a=\pi \)}{\( a=4 \)}{\( a=1 \)}{}{}}
\End

\Data\ID{1003027714}
\Updated{2020-07-15 03:07:12}
\Level{A}
\SubArea{Definite integral}
\LANGen{
\Question{
Which of the given values of \( a \) makes the statement true?
\[ \int\limits_0^1x^8\,\mathrm{d}x = a\cdot\int\limits_{-1}^0x^{24}\,\mathrm{d}x\]}
{\( a=2.\overline{7} \)}{\( a=3 \)}{\( a=0.36 \)}{\( a=16 \)}{}{}}
\LANGpl{
\Question{
Dla jakiej wartości \( a \) zachodzi równość?
\[ \int\limits_0^1x^8\,\mathrm{d}x = a\cdot\int\limits_{-1}^0x^{24}\,\mathrm{d}x\]}
{\( a=2{,}\overline{7} \)}{\( a=3 \)}{\( a=0{,}36 \)}{\( a=16 \)}{}{}}
\LANGcs{
\Question{
Pro kterou z daných hodnot \( a \) platí rovnost?
\[ \int\limits_0^1x^8\,\mathrm{d}x = a\cdot\int\limits_{-1}^0x^{24}\,\mathrm{d}x\]}
{\( a=2{,}\overline{7} \)}{\( a=3 \)}{\( a=0{,}36 \)}{\( a=16 \)}{}{}}
\LANGsk{
\Question{
Pre ktorú z daných hodnôt \( a \) platí rovnosť?
\[ \int\limits_0^1x^8\,\mathrm{d}x = a\cdot\int\limits_{-1}^0x^{24}\,\mathrm{d}x\]}
{\( a=2{,}\overline{7} \)}{\( a=3 \)}{\( a=0{,}36 \)}{\( a=16 \)}{}{}}
\LANGes{
\Question{
¿Cuál de los valores dados de \( a \) hace el enunciado verdadero?
\[ \int\limits_0^1x^8\,\mathrm{d}x = a\cdot\int\limits_{-1}^0x^{24}\,\mathrm{d}x\]}
{\( a=2.\overline{7} \)}{\( a=3 \)}{\( a=0.36 \)}{\( a=16 \)}{}{}}
\End

\Data\ID{1003108001}
\Updated{2020-07-15 03:07:12}
\Level{A}
\SubArea{Definite integral}
\LANGen{
\Question{
Evaluate the definite integral. \[ \int\limits_1^2\left(x^7-\ln7\cdot7^x+7x\right)\mathrm{d}x \]}
{\( \frac38 \)}{\( -\frac{51}8 \)}{\( \frac{51}8 \)}{\( \frac{61}8 \)}{}{}}
\LANGpl{
\Question{
Oblicz całkę oznaczoną. \[ \int\limits_1^2\left(x^7-\ln7\cdot7^x+7x\right)\mathrm{d}x \]}
{\( \frac38 \)}{\( -\frac{51}8 \)}{\( \frac{51}8 \)}{\( \frac{61}8 \)}{}{}}
\LANGcs{
\Question{
Vypočítejte určitý integrál. \[ \int\limits_1^2\left(x^7-\ln7\cdot7^x+7x\right)\mathrm{d}x \]}
{\( \frac38 \)}{\( -\frac{51}8 \)}{\( \frac{51}8 \)}{\( \frac{61}8 \)}{}{}}
\LANGsk{
\Question{
Vypočítajte určitý integrál. \[ \int\limits_1^2\left(x^7-\ln7\cdot7^x+7x\right)\mathrm{d}x \]}
{\( \frac38 \)}{\( -\frac{51}8 \)}{\( \frac{51}8 \)}{\( \frac{61}8 \)}{}{}}
\LANGes{
\Question{
Evalúa la siguiente integral definida. \[ \int\limits_1^2\left(x^7-\ln7\cdot7^x+7x\right)\mathrm{d}x \]}
{\( \frac38 \)}{\( -\frac{51}8 \)}{\( \frac{51}8 \)}{\( \frac{61}8 \)}{}{}}
\End

\Data\ID{1003108002}
\Updated{2020-07-15 03:07:12}
\Level{A}
\SubArea{Definite integral}
\LANGen{
\Question{
Evaluate the definite integral. \[ \int\limits_0^1\left(\frac53\cdot\sqrt[3]{x^2}-x^3+\ln 3\right)\mathrm{d}x \]}
{\( \frac34+\ln 3 \)}{\( \frac54+\ln 3 \)}{\( \frac34 \)}{\( \frac54 \)}{}{}}
\LANGpl{
\Question{
Oblicz całkę oznaczoną. \[ \int\limits_0^1\left(\frac53\cdot\sqrt[3]{x^2}-x^3+\ln 3\right)\mathrm{d}x \]}
{\( \frac34+\ln 3 \)}{\( \frac54+\ln 3 \)}{\( \frac34 \)}{\( \frac54 \)}{}{}}
\LANGcs{
\Question{
Vypočítejte určitý integrál. \[ \int\limits_0^1\left(\frac53\cdot\sqrt[3]{x^2}-x^3+\ln 3\right)\mathrm{d}x \]}
{\( \frac34+\ln 3 \)}{\( \frac54+\ln 3 \)}{\( \frac34 \)}{\( \frac54 \)}{}{}}
\LANGsk{
\Question{
Vypočítajte určitý integrál. \[ \int\limits_0^1\left(\frac53\cdot\sqrt[3]{x^2}-x^3+\ln 3\right)\mathrm{d}x \]}
{\( \frac34+\ln 3 \)}{\( \frac54+\ln 3 \)}{\( \frac34 \)}{\( \frac54 \)}{}{}}
\LANGes{
\Question{
Evalúa la siguiente integral definida. \[ \int\limits_0^1\left(\frac53\cdot\sqrt[3]{x^2}-x^3+\ln 3\right)\mathrm{d}x \]}
{\( \frac34+\ln 3 \)}{\( \frac54+\ln 3 \)}{\( \frac34 \)}{\( \frac54 \)}{}{}}
\End

\Data\ID{1003108003}
\Updated{2020-07-15 03:07:12}
\Level{A}
\SubArea{Definite integral}
\LANGen{
\Question{
Evaluate the  definite integral. \[ \int\limits_{-\frac{\pi}3}^{\frac{\pi}3}(\sin x-\cos x )\mathrm{d}x \]}
{\( -\sqrt3 \)}{\( 1 \)}{\( 1-\sqrt3 \)}{\( 1+\sqrt3 \)}{}{}}
\LANGpl{
\Question{
Oblicz całkę oznaczoną. \[ \int\limits_{-\frac{\pi}3}^{\frac{\pi}3}(\sin x-\cos x )\mathrm{d}x \]}
{\( -\sqrt3 \)}{\( 1 \)}{\( 1-\sqrt3 \)}{\( 1+\sqrt3 \)}{}{}}
\LANGcs{
\Question{
Vypočítejte určitý integrál. \[ \int\limits_{-\frac{\pi}3}^{\frac{\pi}3}(\sin x-\cos x )\mathrm{d}x \]}
{\( -\sqrt3 \)}{\( 1 \)}{\( 1-\sqrt3 \)}{\( 1+\sqrt3 \)}{}{}}
\LANGsk{
\Question{
Vypočítajte určitý integrál. \[ \int\limits_{-\frac{\pi}3}^{\frac{\pi}3}(\sin x-\cos x )\mathrm{d}x \]}
{\( -\sqrt3 \)}{\( 1 \)}{\( 1-\sqrt3 \)}{\( 1+\sqrt3 \)}{}{}}
\LANGes{
\Question{
Evalúa la siguiente integral definida.  \[ \int\limits_{-\frac{\pi}3}^{\frac{\pi}3}(\sin x-\cos x )\mathrm{d}x \]}
{\( -\sqrt3 \)}{\( 1 \)}{\( 1-\sqrt3 \)}{\( 1+\sqrt3 \)}{}{}}
\End

\Data\ID{1003108004}
\Updated{2020-07-15 03:07:12}
\Level{A}
\SubArea{Definite integral}
\LANGen{
\Question{
Evaluate the  definite integral. \[ \int\limits_0^1\left(6\sqrt x-5x^4+3\mathrm{e}^x\right)\mathrm{d}x \]}
{\( 3\mathrm{e} \)}{\( 6+3\mathrm{e} \)}{\( 0 \)}{\( -3\mathrm{e} \)}{}{}}
\LANGpl{
\Question{
Oblicz całkę oznaczoną. \[ \int\limits_0^1\left(6\sqrt x-5x^4+3\mathrm{e}^x\right)\mathrm{d}x \]}
{\( 3\mathrm{e} \)}{\( 6+3\mathrm{e} \)}{\( 0 \)}{\( -3\mathrm{e} \)}{}{}}
\LANGcs{
\Question{
Vypočítejte určitý integrál. \[ \int\limits_0^1\left(6\sqrt x-5x^4+3\mathrm{e}^x\right)\mathrm{d}x \]}
{\( 3\mathrm{e} \)}{\( 6+3\mathrm{e} \)}{\( 0 \)}{\( -3\mathrm{e} \)}{}{}}
\LANGsk{
\Question{
Vypočítajte určitý integrál. \[ \int\limits_0^1\left(6\sqrt x-5x^4+3\mathrm{e}^x\right)\mathrm{d}x \]}
{\( 3\mathrm{e} \)}{\( 6+3\mathrm{e} \)}{\( 0 \)}{\( -3\mathrm{e} \)}{}{}}
\LANGes{
\Question{
Evalúa la siguiente integral definida. \[ \int\limits_0^1\left(6\sqrt x-5x^4+3\mathrm{e}^x\right)\mathrm{d}x \]}
{\( 3\mathrm{e} \)}{\( 6+3\mathrm{e} \)}{\( 0 \)}{\( -3\mathrm{e} \)}{}{}}
\End

\Data\ID{1003108005}
\Updated{2020-07-15 03:07:12}
\Level{A}
\SubArea{Definite integral}
\LANGen{
\Question{
Evaluate the  definite integral. \[ \int\limits_{\frac{\pi}6}^{\frac{\pi}3}\left(\frac3{\cos^2x} -\frac6{\sin^2x}\right)\mathrm{d}x \]}
{\( -2\sqrt3 \)}{\( 2\sqrt3 \)}{\( 0 \)}{\( 12\sqrt3 \)}{}{}}
\LANGpl{
\Question{
Oblicz całkę oznaczoną. \[ \int\limits_{\frac{\pi}6}^{\frac{\pi}3}\left(\frac3{\cos^2x} -\frac6{\sin^2x}\right)\mathrm{d}x \]}
{\( -2\sqrt3 \)}{\( 2\sqrt3 \)}{\( 0 \)}{\( 12\sqrt3 \)}{}{}}
\LANGcs{
\Question{
Vypočítejte určitý integrál. \[ \int\limits_{\frac{\pi}6}^{\frac{\pi}3}\left(\frac3{\cos^2x} -\frac6{\sin^2x}\right)\mathrm{d}x \]}
{\( -2\sqrt3 \)}{\( 2\sqrt3 \)}{\( 0 \)}{\( 12\sqrt3 \)}{}{}}
\LANGsk{
\Question{
Vypočítajte určitý integrál. \[ \int\limits_{\frac{\pi}6}^{\frac{\pi}3}\left(\frac3{\cos^2x} -\frac6{\sin^2x}\right)\mathrm{d}x \]}
{\( -2\sqrt3 \)}{\( 2\sqrt3 \)}{\( 0 \)}{\( 12\sqrt3 \)}{}{}}
\LANGes{
\Question{
Evalúa la siguiente integral definida. \[ \int\limits_{\frac{\pi}6}^{\frac{\pi}3}\left(\frac3{\cos^2x} -\frac6{\sin^2x}\right)\mathrm{d}x \]}
{\( -2\sqrt3 \)}{\( 2\sqrt3 \)}{\( 0 \)}{\( 12\sqrt3 \)}{}{}}
\End

\Data\ID{1003108006}
\Updated{2020-07-15 03:07:12}
\Level{A}
\SubArea{Definite integral}
\LANGen{
\Question{
Evaluate the  definite integral. \[ \int\limits_1^2\left(\frac2x-\frac x2+x^2\right)\mathrm{d}x \]}
{\( \frac{19}{12}+\ln 4 \)}{\( \frac{19}{12}+\ln 2 \)}{\( \frac{43}{12} +\ln 4 \)}{\( \frac{19}{12} \)}{}{}}
\LANGpl{
\Question{
Oblicz całkę oznaczoną. \[ \int\limits_1^2\left(\frac2x-\frac x2+x^2\right)\mathrm{d}x \]}
{\( \frac{19}{12}+\ln 4 \)}{\( \frac{19}{12}+\ln 2 \)}{\( \frac{43}{12} +\ln 4 \)}{\( \frac{19}{12} \)}{}{}}
\LANGcs{
\Question{
Vypočítejte určitý integrál. \[ \int\limits_1^2\left(\frac2x-\frac x2+x^2\right)\mathrm{d}x \]}
{\( \frac{19}{12}+\ln 4 \)}{\( \frac{19}{12}+\ln 2 \)}{\( \frac{43}{12} +\ln 4 \)}{\( \frac{19}{12} \)}{}{}}
\LANGsk{
\Question{
Vypočítajte určitý integrál. \[ \int\limits_1^2\left(\frac2x-\frac x2+x^2\right)\mathrm{d}x \]}
{\( \frac{19}{12}+\ln 4 \)}{\( \frac{19}{12}+\ln 2 \)}{\( \frac{43}{12} +\ln 4 \)}{\( \frac{19}{12} \)}{}{}}
\LANGes{
\Question{
Evalúa la siguiente integral definida.  \[ \int\limits_1^2\left(\frac2x-\frac x2+x^2\right)\mathrm{d}x \]}
{\( \frac{19}{12}+\ln 4 \)}{\( \frac{19}{12}+\ln 2 \)}{\( \frac{43}{12} +\ln 4 \)}{\( \frac{19}{12} \)}{}{}}
\End

\Data\ID{1003108007}
\Updated{2020-07-15 03:07:12}
\Level{A}
\SubArea{Definite integral}
\LANGen{
\Question{
What is the difference \( I_1-I_2 \), where \( I_1=\int\limits_{-1}^1\left(\frac x2+2\right)\mathrm{d}x \) and \( I_2=\int\limits_1^2\frac2x\mathrm{d}x \)?}
{\( 4-\ln 4\)}{\( \ln 4 \)}{\( 4 + \ln 4 \)}{\( \ln 4 - 4 \)}{}{}}
\LANGpl{
\Question{
Wskaż różnicę \( I_1-I_2 \), jeśli \( I_1=\int\limits_{-1}^1\left(\frac x2+2\right)\mathrm{d}x \) i \( I_2=\int\limits_1^2\frac2x\mathrm{d}x \)?}
{\( 4-\ln 4\)}{\( \ln 4 \)}{\( 4 + \ln 4 \)}{\( \ln 4 - 4 \)}{}{}}
\LANGcs{
\Question{
O kolik je větší  \( I_1=\int\limits_{-1}^1\left(\frac x2+2\right)\mathrm{d}x \) než \( I_2=\int\limits_1^2\frac2x\mathrm{d}x \)?}
{o \( 4-\ln 4\)}{o \( \ln 4 \)}{o \( 4 + \ln 4 \)}{o \( \ln 4 - 4 \)}{}{}}
\LANGsk{
\Question{
O koľko je väčší  \( I_1=\int\limits_{-1}^1\left(\frac x2+2\right)\mathrm{d}x \) než \( I_2=\int\limits_1^2\frac2x\mathrm{d}x \)?}
{o \( 4-\ln 4\)}{o \( \ln 4 \)}{o \( 4 + \ln 4 \)}{o \( \ln 4 - 4 \)}{}{}}
\LANGes{
\Question{
¿Cuál es la diferencia \( I_1-I_2 \), donde \( I_1=\int\limits_{-1}^1\left(\frac x2+2\right)\mathrm{d}x \) y \( I_2=\int\limits_1^2\frac2x\mathrm{d}x \)?}
{\( 4-\ln 4\)}{\( \ln 4 \)}{\( 4 + \ln 4 \)}{\( \ln 4 - 4 \)}{}{}}
\End

\Data\ID{1003108008}
\Updated{2020-07-15 03:07:12}
\Level{A}
\SubArea{Definite integral}
\LANGen{
\Question{
Find the value of \( p \) such  that \[ \int\limits_{-2}^3\left(3x^2-5x^4+4x+p\right)\mathrm{d}x=-255. \]}
{\( -5 \)}{\( 5 \)}{\( 3 \)}{No such \( p \) exists.}{}{}}
\LANGpl{
\Question{
Oblicz wartość parametru \( p \) taką, że \[ \int\limits_{-2}^3\left(3x^2-5x^4+4x+p\right)\mathrm{d}x=-255. \]}
{\( -5 \)}{\( 5 \)}{\( 3 \)}{Nie ma takiej wartości \( p \).}{}{}}
\LANGcs{
\Question{
Určete hodnotu parametru \( p \) tak, aby platilo \[ \int\limits_{-2}^3\left(3x^2-5x^4+4x+p\right)\mathrm{d}x=-255. \]}
{\( -5 \)}{\( 5 \)}{\( 3 \)}{Takové \( p \) neexistuje.}{}{}}
\LANGsk{
\Question{
Určte hodnotu parametra \( p \) tak, aby platilo \[ \int\limits_{-2}^3\left(3x^2-5x^4+4x+p\right)\mathrm{d}x=-255 \]}
{\( -5 \)}{\( 5 \)}{\( 3 \)}{Také \( p \) neexistuje.}{}{}}
\LANGes{
\Question{
Halla el valor de \( p \) para el cual \[ \int\limits_{-2}^3\left(3x^2-5x^4+4x+p\right)\mathrm{d}x=-255. \]}
{\( -5 \)}{\( 5 \)}{\( 3 \)}{Este valor de \( p \) no existe.}{}{}}
\End

\Data\ID{1003108106}
\Updated{2020-07-15 03:07:21}
\Level{A}
\SubArea{Definite integral}
\LANGen{
\Question{
How many times is \( \int\limits_0^{\frac{\pi}2}3\cos x\,\mathrm{d}x \) bigger than \( \int\limits_{\frac{\pi}2}^{\pi}\frac{\sin x}2\,\mathrm{d}x \)?}
{\( 6 \) times}{\( 3 \) times}{\( 2 \) times}{It is not bigger.}{}{}}
\LANGpl{
\Question{
Ile razy \( \int\limits_0^{\frac{\pi}2}3\cos x\,\mathrm{d}x \) jest większa niż \( \int\limits_{\frac{\pi}2}^{\pi}\frac{\sin x}2\,\mathrm{d}x \)?}
{\( 6 \) razy}{\( 3 \) razy}{\( 2 \) razy}{To nie jest większe.}{}{}}
\LANGcs{
\Question{
Kolikrát je \( \int\limits_0^{\frac{\pi}2}3\cos x\,\mathrm{d}x \) větší než \( \int\limits_{\frac{\pi}2}^{\pi}\frac{\sin x}2\,\mathrm{d}x \)?}
{\( 6 \) krát}{\( 3 \) krát}{\( 2 \) krát}{Není větší.}{}{}}
\LANGsk{
\Question{
Koľkokrát je \( \int\limits_0^{\frac{\pi}2}3\cos x\,\mathrm{d}x \) väčšia než \( \int\limits_{\frac{\pi}2}^{\pi}\frac{\sin x}2\,\mathrm{d}x \)?}
{\( 6 \) krát}{\( 3 \) krát}{\( 2 \) krát}{Nie je väčšia.}{}{}}
\LANGes{
\Question{
¿Cuántas veces más grande es \( \int\limits_0^{\frac{\pi}2}3\cos x\,\mathrm{d}x \) con respecto a \( \int\limits_{\frac{\pi}2}^{\pi}\frac{\sin x}2\,\mathrm{d}x \)?}
{\( 6 \) veces}{\( 3 \) veces}{\( 2 \) veces}{No es más grande.}{}{}}
\End

\Data\ID{9000150401}
\Updated{2020-07-15 03:07:37}
\Level{A}
\SubArea{Definite integral}
\LANGen{
\Question{
Evaluate the following definite integral. 
\[ 
 \int _{-3}^{1}(x^{2} + 3x)\, \text{d}x
\]}
{\(-\frac{8} 
{3}\)}{\(\frac{8}
{3}\)}{\(-\frac{64} 
 {3} \)}{\(\frac{64}
 {3} \)}{}{}}
\LANGpl{
\Question{
Oblicz.
\[ 
 \int _{-3}^{1}(x^{2} + 3x)\, \text{d}x
\]}
{\(-\frac{8} 
{3}\)}{\(\frac{8}
{3}\)}{\(-\frac{64} 
 {3} \)}{\(\frac{64}
 {3} \)}{}{}}
\LANGcs{
\Question{
Vypočítejte \(\int _{-3}^{1}(x^{2} + 3x)\, \text{d}x\).}
{\(-\frac{8} 
{3}\)}{\(\frac{8}
{3}\)}{\(-\frac{64} 
 {3} \)}{\(\frac{64}
 {3} \)}{}{}}
\LANGsk{
\Question{
Vypočítajte \(\int _{-3}^{1}(x^{2} + 3x)\, \text{d}x\).}
{\(-\frac{8} 
{3}\)}{\(\frac{8}
{3}\)}{\(-\frac{64} 
 {3} \)}{\(\frac{64}
 {3} \)}{}{}}
\LANGes{
\Question{
Evalúa la siguiente integral definida.
\[ 
 \int _{-3}^{1}(x^{2} + 3x)\, \text{d}x
\]}
{\(-\frac{8} 
{3}\)}{\(\frac{8}
{3}\)}{\(-\frac{64} 
 {3} \)}{\(\frac{64}
 {3} \)}{}{}}
\End

\Data\ID{9000150402}
\Updated{2020-07-15 03:07:37}
\Level{A}
\SubArea{Definite integral}
\LANGen{
\Question{
Evaluate the following definite integral. 
\[ 
 \int _{-\frac{\pi }{
2} }^{ \frac{\pi } 
{2} }\sin x\, \text{d}x
\]}
{\(0\)}{\(\pi \)}{\(2\)}{\(1\)}{}{}}
\LANGpl{
\Question{
Oblicz.
\[ 
 \int _{-\frac{\pi }{
2} }^{ \frac{\pi } 
{2} }\sin x\, \text{d}x
\]}
{\(0\)}{\(\pi \)}{\(2\)}{\(1\)}{}{}}
\LANGcs{
\Question{
Vypočítejte \(\int _{-\frac{\pi }{
2} }^{ \frac{\pi } 
{2} }\sin x\, \text{d}x\).}
{\(0\)}{\(\pi \)}{\(2\)}{\(1\)}{}{}}
\LANGsk{
\Question{
Vypočítajte \(\int _{-\frac{\pi }{
2} }^{ \frac{\pi } 
{2} }\sin x\, \text{d}x\).}
{\(0\)}{\(\pi \)}{\(2\)}{\(1\)}{}{}}
\LANGes{
\Question{
Evalúa la siguiente integral definida.
\[ 
 \int _{-\frac{\pi }{
2} }^{ \frac{\pi } 
{2} }\sin x\, \text{d}x
\]}
{\(0\)}{\(\pi \)}{\(2\)}{\(1\)}{}{}}
\End

\Data\ID{9000150403}
\Updated{2020-07-15 03:07:37}
\Level{A}
\SubArea{Definite integral}
\LANGen{
\Question{
Evaluate the following definite integral. 
\[ 
 \int _{-2}^{0}\mathrm{e}^{x}\, \text{d}x
\]}
{\(1 -\frac{1} 
{\mathrm{e}^{2}} \)}{\(1 + \frac{1} 
{\mathrm{e}^{2}} \)}{\(\frac{1}
{\mathrm{e}^{2}} \)}{\(-\frac{1} 
{\mathrm{e}^{2}} \)}{}{}}
\LANGpl{
\Question{
Oblicz.
\[ 
 \int _{-2}^{0}\mathrm{e}^{x}\, \text{d}x
\]}
{\(1 -\frac{1} 
{\mathrm{e}^{2}} \)}{\(1 + \frac{1} 
{\mathrm{e}^{2}} \)}{\(\frac{1}
{\mathrm{e}^{2}} \)}{\(-\frac{1} 
{\mathrm{e}^{2}} \)}{}{}}
\LANGcs{
\Question{
Vypočítejte \(\int _{-2}^{0}\mathrm{e}^{x}\, \text{d}x\).}
{\(1 -\frac{1} 
{\mathrm{e}^{2}} \)}{\(1 + \frac{1} 
{\mathrm{e}^{2}} \)}{\(\frac{1}
{\mathrm{e}^{2}} \)}{\(-\frac{1} 
{\mathrm{e}^{2}} \)}{}{}}
\LANGsk{
\Question{
Vypočítajte \(\int _{-2}^{0}\mathrm{e}^{x}\, \text{d}x\).}
{\(1 -\frac{1} 
{\mathrm{e}^{2}} \)}{\(1 + \frac{1} 
{\mathrm{e}^{2}} \)}{\(\frac{1}
{\mathrm{e}^{2}} \)}{\(-\frac{1} 
{\mathrm{e}^{2}} \)}{}{}}
\LANGes{
\Question{
Evalúa la siguiente integral definida.
\[ 
 \int _{-2}^{0}\mathrm{e}^{x}\, \text{d}x
\]}
{\(1 -\frac{1} 
{\mathrm{e}^{2}} \)}{\(1 + \frac{1} 
{\mathrm{e}^{2}} \)}{\(\frac{1}
{\mathrm{e}^{2}} \)}{\(-\frac{1} 
{\mathrm{e}^{2}} \)}{}{}}
\End

\Data\ID{9000150404}
\Updated{2020-07-15 03:07:37}
\Level{A}
\SubArea{Definite integral}
\LANGen{
\Question{
Evaluate the following definite integral. 
\[ 
 \int _{2}^{6} \frac{2}
{x}\, \text{d}x
\]}
{\(\ln 9\)}{\(\ln 2\)}{\(\ln 3\)}{\(2\ln 6\)}{}{}}
\LANGpl{
\Question{
Oblicz.
\[ 
 \int _{2}^{6} \frac{2}
{x}\, \text{d}x
\]}
{\(\ln 9\)}{\(\ln 2\)}{\(\ln 3\)}{\(2\ln 6\)}{}{}}
\LANGcs{
\Question{
Vypočítejte \(\int _{2}^{6} \frac{2}
{x}\, \text{d}x\).}
{\(\ln 9\)}{\(\ln 2\)}{\(\ln 3\)}{\(2\ln 6\)}{}{}}
\LANGsk{
\Question{
Vypočítajte \(\int _{2}^{6} \frac{2}
{x}\, \text{d}x\).}
{\(\ln 9\)}{\(\ln 2\)}{\(\ln 3\)}{\(2\ln 6\)}{}{}}
\LANGes{
\Question{
Evalúa la siguiente integral definida.
\[ 
 \int _{2}^{6} \frac{2}
{x}\, \text{d}x
\]}
{\(\ln 9\)}{\(\ln 2\)}{\(\ln 3\)}{\(2\ln 6\)}{}{}}
\End

\Data\ID{9000150407}
\Updated{2020-07-15 03:07:37}
\Level{A}
\SubArea{Definite integral}
\LANGen{
\Question{
Evaluate the following definite integral. 
\[ 
 \int _{1}^{2}7^{x}\, \text{d}x
\]}
{\(\frac{42}
 {\ln 7} \)}{\(49\ln 7\)}{\(42\)}{\(42\ln 7\)}{}{}}
\LANGpl{
\Question{
Oblicz.
\[ 
 \int _{1}^{2}7^{x}\, \text{d}x
\]}
{\(\frac{42}
 {\ln 7} \)}{\(49\ln 7\)}{\(42\)}{\(42\ln 7\)}{}{}}
\LANGcs{
\Question{
Vypočítejte \(\int _{1}^{2}7^{x}\, \text{d}x\).}
{\(\frac{42}
 {\ln 7} \)}{\(49\ln 7\)}{\(42\)}{\(42\ln 7\)}{}{}}
\LANGsk{
\Question{
Vypočítajte \(\int _{1}^{2}7^{x}\, \text{d}x\).}
{\(\frac{42}
 {\ln 7} \)}{\(49\ln 7\)}{\(42\)}{\(42\ln 7\)}{}{}}
\LANGes{
\Question{
Evalúa la siguiente integral definida.
\[ 
 \int _{1}^{2}7^{x}\, \text{d}x
\]}
{\(\frac{42}
 {\ln 7} \)}{\(49\ln 7\)}{\(42\)}{\(42\ln 7\)}{}{}}
\End

\Data\ID{9000150408}
\Updated{2020-07-15 03:07:37}
\Level{A}
\SubArea{Definite integral}
\LANGen{
\Question{
Evaluate the following definite integral. 
\[ 
 \int _{0}^{ \frac{\pi }{
4} } \frac{2}
{\cos ^{2}x}\, \text{d}x
\]}
{\(2\)}{\(0\)}{\(4\)}{\(\pi \)}{}{}}
\LANGpl{
\Question{
Oblicz.
\[ 
 \int _{0}^{ \frac{\pi }{
4} } \frac{2}
{\cos ^{2}x}\, \text{d}x
\]}
{\(2\)}{\(0\)}{\(4\)}{\(\pi \)}{}{}}
\LANGcs{
\Question{
Vypočítejte \(\int _{0}^{ \frac{\pi }{
4} } \frac{2}
{\cos ^{2}x}\, \text{d}x\).}
{\(2\)}{\(0\)}{\(4\)}{\(\pi \)}{}{}}
\LANGsk{
\Question{
Vypočítajte \(\int _{0}^{ \frac{\pi }{
4} } \frac{2}
{\cos ^{2}x}\, \text{d}x\).}
{\(2\)}{\(0\)}{\(4\)}{\(\pi \)}{}{}}
\LANGes{
\Question{
Evalúa la siguiente integral definida.
\[ 
 \int _{0}^{ \frac{\pi }{
4} } \frac{2}
{\cos ^{2}x}\, \text{d}x
\]}
{\(2\)}{\(0\)}{\(4\)}{\(\pi \)}{}{}}
\End

\Data\ID{1003027601}
\Updated{2020-07-15 03:07:12}
\Level{B}
\SubArea{Definite integral}
\LANGen{
\Question{
How much bigger is \( \int\limits_0^3\left[(x-3)^2+1\right]\!\mathrm{d}x \) than \( \int\limits_3^6 \log_55\,\mathrm{d}x \)?}
{\( 9 \)}{\( 15 \)}{\( 4 \)}{none of the options}{}{}}
\LANGpl{
\Question{
O ile zwiększa się wartość \( \int\limits_0^3\left[(x-3)^2+1\right]\!\mathrm{d}x \) od \( \int\limits_3^6 \log_55\,\mathrm{d}x \)?}
{\( 9 \)}{\( 15 \)}{\( 4 \)}{żadna z opcji}{}{}}
\LANGcs{
\Question{
O kolik je větší \( \int\limits_0^3\left[(x-3)^2+1\right]\!\mathrm{d}x \) než \( \int\limits_3^6 \log_55\,\mathrm{d}x \)?}
{o \( 9 \)}{o \( 15 \)}{o \( 4 \)}{žádná z uvedených možností není správná}{}{}}
\LANGsk{
\Question{
O koľko je väčší \( \int\limits_0^3\left[(x-3)^2+1\right]\!\mathrm{d}x \) než \( \int\limits_3^6 \log_55\,\mathrm{d}x \)?}
{o \( 9 \)}{o \( 15 \)}{o \( 4 \)}{žiadna z uvedených možností nie je správna}{}{}}
\LANGes{
\Question{
¿Cuántas veces más grande es \( \int\limits_0^3\left[(x-3)^2+1\right]\!\mathrm{d}x \) con respecto a \( \int\limits_3^6 \log_55\,\mathrm{d}x \)?}
{\( 9 \)}{\( 15 \)}{\( 4 \)}{ninguna de las posibilidades}{}{}}
\End

\Data\ID{1003027602}
\Updated{2020-07-15 03:07:12}
\Level{B}
\SubArea{Definite integral}
\LANGen{
\Question{
How much smaller is \( \int\limits_1^{10}\frac{x+1}{x^2+x}\,\mathrm{d}x \) than \( \int\limits_1^{10}\frac{11-x}{10}\,\mathrm{d}x \)?}
{more than \( 1 \)}{less than \( 1 \)}{they are equal}{cannot be determined}{}{}}
\LANGpl{
\Question{
O ile mniejsza jest  \( \int\limits_1^{10}\frac{x+1}{x^2+x}\,\mathrm{d}x \) od \( \int\limits_1^{10}\frac{11-x}{10}\,\mathrm{d}x \)?}
{więcej niż \( 1 \)}{mniej niż \( 1 \)}{są równe}{nie można tego określić}{}{}}
\LANGcs{
\Question{
O kolik je menší \( \int\limits_1^{10}\frac{x+1}{x^2+x}\,\mathrm{d}x \) než \( \int\limits_1^{10}\frac{11-x}{10}\,\mathrm{d}x \)?}
{o více než \( 1 \)}{o méně než \( 1 \)}{jsou stejné}{nelze určit}{}{}}
\LANGsk{
\Question{
O koľko je menší \( \int\limits_1^{10}\frac{x+1}{x^2+x}\,\mathrm{d}x \) než \( \int\limits_1^{10}\frac{11-x}{10}\,\mathrm{d}x \)?}
{o viac ako \( 1 \)}{o menej ako \( 1 \)}{sú rovnaké}{nedá sa určiť}{}{}}
\LANGes{
\Question{
¿Cuánto más pequeña es la integral \( \int\limits_1^{10}\frac{x+1}{x^2+x}\,\mathrm{d}x \) con respecto a \( \int\limits_1^{10}\frac{11-x}{10}\,\mathrm{d}x \)?}
{más de \( 1 \)}{menos de \( 1 \)}{son iguales}{no se puede determinar}{}{}}
\End

\Data\ID{1003027603}
\Updated{2020-07-15 03:07:12}
\Level{B}
\SubArea{Definite integral}
\LANGen{
\Question{
How many times is \( \int\limits_{\frac{\pi}6}^{\frac{\pi}3} \frac{\sin 2x}{\cos x}\,\mathrm{d}x \) bigger than \( \int\limits_{-\frac{\pi}3}^{-\frac{\pi}6} \frac{\sin 2x}{\sin x}\,\mathrm{d}x \)?}
{they are equal}{\( 2 \) times}{\( 4 \) times}{none of the options}{}{}}
\LANGpl{
\Question{
Ile razy jest  \( \int\limits_{\frac{\pi}6}^{\frac{\pi}3} \frac{\sin 2x}{\cos x}\,\mathrm{d}x \) większa od \( \int\limits_{-\frac{\pi}3}^{-\frac{\pi}6} \frac{\sin 2x}{\sin x}\,\mathrm{d}x \)?}
{są równe}{\( 2 \) razy}{\( 4 \) razy}{żadna z opcji}{}{}}
\LANGcs{
\Question{
Kolikrát je \( \int\limits_{\frac{\pi}6}^{\frac{\pi}3} \frac{\sin 2x}{\cos x}\,\mathrm{d}x \) větší než \( \int\limits_{-\frac{\pi}3}^{-\frac{\pi}6} \frac{\sin 2x}{\sin x}\,\mathrm{d}x \)?}
{jsou stejné}{\( 2 \) krát}{\( 4 \) krát}{žádná z uvedených možností není správná}{}{}}
\LANGsk{
\Question{
Koľkokrát je \( \int\limits_{\frac{\pi}6}^{\frac{\pi}3} \frac{\sin 2x}{\cos x}\,\mathrm{d}x \) väčší než \( \int\limits_{-\frac{\pi}3}^{-\frac{\pi}6} \frac{\sin 2x}{\sin x}\,\mathrm{d}x \)?}
{sú rovnaké}{\( 2 \) krát}{\( 4 \) krát}{žiadna z uvedených možností nie je správna}{}{}}
\LANGes{
\Question{
¿Cuántas veces más grande es \( \int\limits_{\frac{\pi}6}^{\frac{\pi}3} \frac{\sin 2x}{\cos x}\,\mathrm{d}x \) con respecto a \( \int\limits_{-\frac{\pi}3}^{-\frac{\pi}6} \frac{\sin 2x}{\sin x}\,\mathrm{d}x \)?}
{son iguales}{\( 2 \) veces}{\( 4 \) veces}{ninguna de las respuestas}{}{}}
\End

\Data\ID{1003027604}
\Updated{2020-07-15 03:07:12}
\Level{B}
\SubArea{Definite integral}
\LANGen{
\Question{
How many times is \( \int\limits_{-3}^0 \frac{x^2+x-6}{x-2}\,\mathrm{d}x \) smaller than \( \int\limits_{4}^7 \frac{5x^2-15x-20}{x+1}\,\mathrm{d}x \)?}
{\( 5 \) times}{\( 9 \) times}{they are equal}{none of the options}{}{}}
\LANGpl{
\Question{
Ile razy \( \int\limits_{-3}^0 \frac{x^2+x-6}{x-2}\,\mathrm{d}x \) jest mniejsza niż  \( \int\limits_{4}^7 \frac{5x^2-15x-20}{x+1}\,\mathrm{d}x \)?}
{\( 5 \) razy}{\( 9 \) razy}{są równe}{żadna z opcji}{}{}}
\LANGcs{
\Question{
Kolikrát je \( \int\limits_{-3}^0 \frac{x^2+x-6}{x-2}\,\mathrm{d}x \) menší než \( \int\limits_{4}^7 \frac{5x^2-15x-20}{x+1}\,\mathrm{d}x \)?}
{\( 5 \) krát}{\( 9 \) krát}{jsou stejné}{žádná z uvedených možností není správná}{}{}}
\LANGsk{
\Question{
Koľkokrát je \( \int\limits_{-3}^0 \frac{x^2+x-6}{x-2}\,\mathrm{d}x \) menší než \( \int\limits_{4}^7 \frac{5x^2-15x-20}{x+1}\,\mathrm{d}x \)?}
{\( 5 \) krát}{\( 9 \) krát}{sú rovnaké}{žiadna z uvedených možností nie je správna}{}{}}
\LANGes{
\Question{
¿Cuántas veces más pequeña es \( \int\limits_{-3}^0 \frac{x^2+x-6}{x-2}\,\mathrm{d}x \) con respecto a \( \int\limits_{4}^7 \frac{5x^2-15x-20}{x+1}\,\mathrm{d}x \)?}
{\( 5 \) veces}{\( 9 \) veces}{son iguales}{ninguna de las respuestas}{}{}}
\End

\Data\ID{1003108101}
\Updated{2020-07-15 03:07:21}
\Level{B}
\SubArea{Definite integral}
\LANGen{
\Question{
Evaluate the integral \( \int\limits_1^4\frac{\sqrt x-x+x^2}x\,\mathrm{d}x \).}
{\( 6.5 \)}{\( 5.5 \)}{\( 9.5 \)}{\( 14.5 \)}{}{}}
\LANGpl{
\Question{
Wyznacz całkę \( \int\limits_1^4\frac{\sqrt x-x+x^2}x\,\mathrm{d}x \).}
{\( 6{,}5 \)}{\( 5{,}5 \)}{\( 9{,}5 \)}{\( 14{,}5 \)}{}{}}
\LANGcs{
\Question{
Vypočítejte integrál \( \int\limits_1^4\frac{\sqrt x-x+x^2}x\,\mathrm{d}x \).}
{\( 6{,}5 \)}{\( 5{,}5 \)}{\( 9{,}5 \)}{\( 14{,}5 \)}{}{}}
\LANGsk{
\Question{
Vypočítajte integrál \( \int\limits_1^4\frac{\sqrt x-x+x^2}x\,\mathrm{d}x \).}
{\( 6{,}5 \)}{\( 5{,}5 \)}{\( 9{,}5 \)}{\( 14{,}5 \)}{}{}}
\LANGes{
\Question{
Evalúa la integral  \( \int\limits_1^4\frac{\sqrt x-x+x^2}x\,\mathrm{d}x \).}
{\( 6.5 \)}{\( 5.5 \)}{\( 9.5 \)}{\( 14.5 \)}{}{}}
\End

\Data\ID{1003108102}
\Updated{2020-07-15 03:07:21}
\Level{B}
\SubArea{Definite integral}
\LANGen{
\Question{
Evaluate the integral \( \int\limits_{-1}^1 \left(x^2-x\right)(x+1)\,\mathrm{d}x \).}
{\( 0 \)}{\( \frac12 \)}{\( 2 \)}{\( 1 \)}{}{}}
\LANGpl{
\Question{
Wyznacz całkę \( \int\limits_{-1}^1 \left(x^2-x\right)(x+1)\,\mathrm{d}x \).}
{\( 0 \)}{\( \frac12 \)}{\( 2 \)}{\( 1 \)}{}{}}
\LANGcs{
\Question{
Vypočítejte integrál \( \int\limits_{-1}^1 \left(x^2-x\right)(x+1)\,\mathrm{d}x \).}
{\( 0 \)}{\( \frac12 \)}{\( 2 \)}{\( 1 \)}{}{}}
\LANGsk{
\Question{
Vypočítajte integrál \( \int\limits_{-1}^1 \left(x^2-x\right)(x+1)\,\mathrm{d}x \).}
{\( 0 \)}{\( \frac12 \)}{\( 2 \)}{\( 1 \)}{}{}}
\LANGes{
\Question{
Evalúa la  integral \( \int\limits_{-1}^1 \left(x^2-x\right)(x+1)\,\mathrm{d}x \).}
{\( 0 \)}{\( \frac12 \)}{\( 2 \)}{\( 1 \)}{}{}}
\End

\Data\ID{1003108103}
\Updated{2020-07-15 03:07:21}
\Level{B}
\SubArea{Definite integral}
\LANGen{
\Question{
Evaluate the integral \( \int\limits_{-1}^1\frac{x^2+3x-10}{x-2}\,\mathrm{d}x \).}
{\( 10 \)}{\( 6 \)}{\( 8 \)}{\( 12 \)}{}{}}
\LANGpl{
\Question{
Wyznacz całkę \( \int\limits_{-1}^1\frac{x^2+3x-10}{x-2}\,\mathrm{d}x \).}
{\( 10 \)}{\( 6 \)}{\( 8 \)}{\( 12 \)}{}{}}
\LANGcs{
\Question{
Vypočítejte integrál \( \int\limits_{-1}^1\frac{x^2+3x-10}{x-2}\,\mathrm{d}x \).}
{\( 10 \)}{\( 6 \)}{\( 8 \)}{\( 12 \)}{}{}}
\LANGsk{
\Question{
Vypočítajte integrál \( \int\limits_{-1}^1\frac{x^2+3x-10}{x-2}\,\mathrm{d}x \).}
{\( 10 \)}{\( 6 \)}{\( 8 \)}{\( 12 \)}{}{}}
\LANGes{
\Question{
Evalúa la integral \( \int\limits_{-1}^1\frac{x^2+3x-10}{x-2}\,\mathrm{d}x \).}
{\( 10 \)}{\( 6 \)}{\( 8 \)}{\( 12 \)}{}{}}
\End

\Data\ID{1003108104}
\Updated{2020-07-15 03:07:21}
\Level{B}
\SubArea{Definite integral}
\LANGen{
\Question{
Evaluate the integral \( \int\limits_{-2}^0 (x+2)^3\,\mathrm{d}x \).}
{\( 4 \)}{\( 28 \)}{\( -28 \)}{\( 12 \)}{}{}}
\LANGpl{
\Question{
Wyznacz całkę \( \int\limits_{-2}^0 (x+2)^3\,\mathrm{d}x \).}
{\( 4 \)}{\( 28 \)}{\( -28 \)}{\( 12 \)}{}{}}
\LANGcs{
\Question{
Vypočítejte integrál \( \int\limits_{-2}^0 (x+2)^3\,\mathrm{d}x \).}
{\( 4 \)}{\( 28 \)}{\( -28 \)}{\( 12 \)}{}{}}
\LANGsk{
\Question{
Vypočítajte integrál \( \int\limits_{-2}^0 (x+2)^3\,\mathrm{d}x \).}
{\( 4 \)}{\( 28 \)}{\( -28 \)}{\( 12 \)}{}{}}
\LANGes{
\Question{
Evalúa la integral \( \int\limits_{-2}^0 (x+2)^3\,\mathrm{d}x \).}
{\( 4 \)}{\( 28 \)}{\( -28 \)}{\( 12 \)}{}{}}
\End

\Data\ID{1003108105}
\Updated{2020-07-15 03:07:21}
\Level{B}
\SubArea{Definite integral}
\LANGen{
\Question{
Evaluate the integral \( \int\limits_0^1\frac{x-1}{x+3}\,\mathrm{d}x \).}
{\( 1+\ln\left(\frac34\right)^4 \)}{\( 1-4\ln12 \)}{\( 4\ln0.75 \)}{\( 4\ln12 \)}{}{}}
\LANGpl{
\Question{
Wyznacz całkę \( \int\limits_0^1\frac{x-1}{x+3}\,\mathrm{d}x \).}
{\( 1+\ln\left(\frac34\right)^4 \)}{\( 1-4\ln12 \)}{\( 4\ln0{,}75 \)}{\( 4\ln12 \)}{}{}}
\LANGcs{
\Question{
Vypočítejte integrál \( \int\limits_0^1\frac{x-1}{x+3}\,\mathrm{d}x \).}
{\( 1+\ln\left(\frac34\right)^4 \)}{\( 1-4\ln12 \)}{\( 4\ln0{,}75 \)}{\( 4\ln12 \)}{}{}}
\LANGsk{
\Question{
Vypočítajte integrál \( \int\limits_0^1\frac{x-1}{x+3}\,\mathrm{d}x \).}
{\( 1+\ln\left(\frac34\right)^4 \)}{\( 1-4\ln12 \)}{\( 4\ln0{,}75 \)}{\( 4\ln12 \)}{}{}}
\LANGes{
\Question{
Evalúa la integral \( \int\limits_0^1\frac{x-1}{x+3}\,\mathrm{d}x \).}
{\( 1+\ln\left(\frac34\right)^4 \)}{\( 1-4\ln12 \)}{\( 4\ln0.75 \)}{\( 4\ln12 \)}{}{}}
\End

\Data\ID{1003108107}
\Updated{2020-07-15 03:07:21}
\Level{B}
\SubArea{Definite integral}
\LANGen{
\Question{
Compare the value of \( \int\limits_1^2\frac{x^2\cdot\sqrt[3]x}{\sqrt[4]{x^3}}\,\mathrm{d}x \) to zero.}
{It is bigger than \( 0 \).}{It is smaller than \( 0 \).}{It equals \( 0 \).}{It cannot be compared.}{}{}}
\LANGpl{
\Question{
Porównaj wartość  \( \int\limits_1^2\frac{x^2\cdot\sqrt[3]x}{\sqrt[4]{x^3}}\,\mathrm{d}x \) do zera.}
{Jest większa od  \( 0 \).}{Jest mniejsza niż zero \( 0 \).}{Jest równa \( 0 \).}{Nie można wykonać porównania.}{}{}}
\LANGcs{
\Question{
Porovnejte s číslem 0 hodnotu \( \int\limits_1^2\frac{x^2\cdot\sqrt[3]x}{\sqrt[4]{x^3}}\,\mathrm{d}x \).}
{Je větší než \( 0 \).}{Je menší než \( 0 \).}{Je rovno \( 0 \).}{Nelze porovnat.}{}{}}
\LANGsk{
\Question{
Porovnajte s číslom 0 hodnotu \( \int\limits_1^2\frac{x^2\cdot\sqrt[3]x}{\sqrt[4]{x^3}}\,\mathrm{d}x \).}
{Je väčšia než \( 0 \).}{Je menšia než \( 0 \).}{Je rovná \( 0 \).}{Nedá sa porovnať.}{}{}}
\LANGes{
\Question{
Compara el valor de \( \int\limits_1^2\frac{x^2\cdot\sqrt[3]x}{\sqrt[4]{x^3}}\,\mathrm{d}x \) con el cero.}
{Es más grande de \( 0 \).}{Es más pequeña de \( 0 \).}{Es igual a \( 0 \).}{No se puede comparar.}{}{}}
\End

\Data\ID{1003108108}
\Updated{2020-07-15 03:07:21}
\Level{B}
\SubArea{Definite integral}
\LANGen{
\Question{
Compare the value of \( \int\limits_1^2\frac{x^2-x}{\sqrt x}\,\mathrm{d}x \) to the number \( \frac4{15} \).}
{It is bigger than \( \frac4{15} \).}{It is smaller than \( \frac4{15} \).}{It equals \( \frac4{15} \).}{It cannot be compared.}{}{}}
\LANGpl{
\Question{
Porównaj wartość \( \int\limits_1^2\frac{x^2-x}{\sqrt x}\,\mathrm{d}x \) z liczbą \( \frac4{15} \).}
{Jest większa niż \( \frac4{15} \).}{Jest mniejsza niż \( \frac4{15} \).}{Jest równa \( \frac4{15} \).}{Nie można wykonać porównania.}{}{}}
\LANGcs{
\Question{
Porovnejte hodnotu \( \int\limits_1^2\frac{x^2-x}{\sqrt x}\,\mathrm{d}x \) s číslem \( \frac4{15} \).}
{Je větší než \( \frac4{15} \).}{Je menší než \( \frac4{15} \).}{Je rovno \( \frac4{15} \).}{Nelze porovnat.}{}{}}
\LANGsk{
\Question{
Porovnajte hodnotu \( \int\limits_1^2\frac{x^2-x}{\sqrt x}\,\mathrm{d}x \) s číslom \( \frac4{15} \).}
{Je väčšia než \( \frac4{15} \).}{Je menšia než \( \frac4{15} \).}{Je rovná \( \frac4{15} \).}{Nedá sa porovnať.}{}{}}
\LANGes{
\Question{
Compara el valor de \( \int\limits_1^2\frac{x^2-x}{\sqrt x}\,\mathrm{d}x \) con el número \( \frac4{15} \).}
{Es más grande de \( \frac4{15} \).}{Es más pequeña de \( \frac4{15} \).}{Es igual a \( \frac4{15} \).}{No se puede comparar.}{}{}}
\End

\Data\ID{1003108201}
\Updated{2020-07-15 03:07:21}
\Level{B}
\SubArea{Definite integral}
\LANGen{
\Question{
Evaluate the  definite integral \( \int\limits_0^{\frac{\pi}6}\frac{3\cos2t}{\cos t+\sin t}\,\mathrm{d}t \).  Which of the following intervals contains the value of the integral?}
{\( (0.8;1.2) \)}{\( (0.4;0.8) \)}{\( (-0.8;-0.1) \)}{\( (-0.1;0.4) \)}{}{}}
\LANGpl{
\Question{
Wyznacz całkę oznaczoną \( \int\limits_0^{\frac{\pi}6}\frac{3\cos2t}{\cos t+\sin t}\,\mathrm{d}t \).  Który z poniższych przedziałów zawiera wartość całki?}
{\( (0{,}8;1{,}2) \)}{\( (0{,}4;0{,}8) \)}{\( (-0{,}8;-0{,}1) \)}{\( (-0{,}1;0{,}4) \)}{}{}}
\LANGcs{
\Question{
Vypočítejte určitý integrál \( \int\limits_0^{\frac{\pi}6}\frac{3\cos2t}{\cos t+\sin t}\,\mathrm{d}t \).  Do kterého z uvedených intervalů patří vypočítaná hodnota?}
{\( (0{,}8;1{,}2) \)}{\( (0{,}4;0{,}8) \)}{\( (-0{,}8;-0{,}1) \)}{\( (-0{,}1;0{,}4) \)}{}{}}
\LANGsk{
\Question{
Vypočítajte určitý integrál \( \int\limits_0^{\frac{\pi}6}\frac{3\cos2t}{\cos t+\sin t}\,\mathrm{d}t \).  Do ktorého z uvedených intervalov patrí vypočítaná hodnota?}
{\( (0{,}8;1{,}2) \)}{\( (0{,}4;0{,}8) \)}{\( (-0{,}8;-0{,}1) \)}{\( (-0{,}1;0{,}4) \)}{}{}}
\LANGes{
\Question{
Evalúa la integral definida \( \int\limits_0^{\frac{\pi}6}\frac{3\cos2t}{\cos t+\sin t}\,\mathrm{d}t \).  ¿Cuál de los siguientes intervalos contiene el valor de la integral?}
{\( (0.8;1.2) \)}{\( (0.4;0.8) \)}{\( (-0.8;-0.1) \)}{\( (-0.1;0.4) \)}{}{}}
\End

\Data\ID{1003108202}
\Updated{2020-07-15 03:07:21}
\Level{B}
\SubArea{Definite integral}
\LANGen{
\Question{
Evaluate the  definite integral \( \int\limits_{\frac{\pi}6}^{\frac{\pi}3}\frac{\mathrm{tg}\,a}{\sin2a}\,\mathrm{d}a \).}
{\( \frac{\sqrt3}3 \)}{\( \frac{2\sqrt3}3 \)}{\( -\frac{2\sqrt3}3 \)}{\( -\frac{\sqrt3}3 \)}{}{}}
\LANGpl{
\Question{
Wyznacz całkę oznaczoną \( \int\limits_{\frac{\pi}6}^{\frac{\pi}3}\frac{\mathrm{tg}\,a}{\sin2a}\,\mathrm{d}a \).}
{\( \frac{\sqrt3}3 \)}{\( \frac{2\sqrt3}3 \)}{\( -\frac{2\sqrt3}3 \)}{\( -\frac{\sqrt3}3 \)}{}{}}
\LANGcs{
\Question{
Vypočítejte určitý integrál \( \int\limits_{\frac{\pi}6}^{\frac{\pi}3}\frac{\mathrm{tg}\,a}{\sin2a}\,\mathrm{d}a \).}
{\( \frac{\sqrt3}3 \)}{\( \frac{2\sqrt3}3 \)}{\( -\frac{2\sqrt3}3 \)}{\( -\frac{\sqrt3}3 \)}{}{}}
\LANGsk{
\Question{
Vypočítajte určitý integrál \( \int\limits_{\frac{\pi}6}^{\frac{\pi}3}\frac{\mathrm{tg}\,a}{\sin2a}\,\mathrm{d}a \).}
{\( \frac{\sqrt3}3 \)}{\( \frac{2\sqrt3}3 \)}{\( -\frac{2\sqrt3}3 \)}{\( -\frac{\sqrt3}3 \)}{}{}}
\LANGes{
\Question{
Evalúa la integral definida. \( \int\limits_{\frac{\pi}6}^{\frac{\pi}3}\frac{\mathrm{tg}\,a}{\sin2a}\,\mathrm{d}a \).}
{\( \frac{\sqrt3}3 \)}{\( \frac{2\sqrt3}3 \)}{\( -\frac{2\sqrt3}3 \)}{\( -\frac{\sqrt3}3 \)}{}{}}
\End

\Data\ID{1003108203}
\Updated{2020-07-15 03:07:21}
\Level{B}
\SubArea{Definite integral}
\LANGen{
\Question{
Compare the definite integral \( I=\int\limits_0^{\frac{\pi}4}\frac{\cos2b}{\cos^2b}\,\mathrm{d}b \) to the number \( \frac{\pi}2 \).}
{\( I \) is smaller than \( \frac{\pi}2 \) by \( 1 \).}{\( I \) is bigger than \( \frac{\pi}2 \) by \( 1 \).}{\( I \) is equal to \( \frac{\pi}2 \).}{\( I \) is smaller than \( \frac{\pi}2 \) by \( \frac{\pi}4 \).}{}{}}
\LANGpl{
\Question{
Porównaj całkę oznaczoną \( I=\int\limits_0^{\frac{\pi}4}\frac{\cos2b}{\cos^2b}\,\mathrm{d}b \) z liczbą \( \frac{\pi}2 \).}
{\( I \) jest mniejsza niż \( \frac{\pi}2 \) o \( 1 \).}{\( I \) jest większa niż \( \frac{\pi}2 \) o \( 1 \).}{\( I \) jest równa \( \frac{\pi}2 \).}{\( I \) jest mniejsza niż \( \frac{\pi}2 \) o \( \frac{\pi}4 \).}{}{}}
\LANGcs{
\Question{
Porovnejte hodnotu určitého integrálu \( I=\int\limits_0^{\frac{\pi}4}\frac{\cos2b}{\cos^2b}\,\mathrm{d}b \) s číslem  \( \frac{\pi}2 \).}
{\( I \) je menší než \( \frac{\pi}2 \) o \( 1 \).}{\( I \) je větší než \( \frac{\pi}2 \) o \( 1 \).}{\( I \) je rovno \( \frac{\pi}2 \).}{\( I \) je menší než \( \frac{\pi}2 \) o \( \frac{\pi}4 \).}{}{}}
\LANGsk{
\Question{
Porovnajte hodnotu určitého integrálu \( I=\int\limits_0^{\frac{\pi}4}\frac{\cos2b}{\cos^2b}\,\mathrm{d}b \) s číslom \( \frac{\pi}2 \).}
{\( I \) je menšia než \( \frac{\pi}2 \) o \( 1 \).}{\( I \) je väčšia než \( \frac{\pi}2 \) o \( 1 \).}{\( I \) je rovná \( \frac{\pi}2 \).}{\( I \) je menšia než \( \frac{\pi}2 \) o \( \frac{\pi}4 \).}{}{}}
\LANGes{
\Question{
Compara la integral definida \( I=\int\limits_0^{\frac{\pi}4}\frac{\cos2b}{\cos^2b}\,\mathrm{d}b \) con el número \( \frac{\pi}2 \).}
{\( I \) es más pequeña de \( \frac{\pi}2 \) de \( 1 \).}{\( I \) es más grande de \( \frac{\pi}2 \) de \( 1 \).}{\( I \) es igual a \( \frac{\pi}2 \).}{\( I \) es más pequeña de \( \frac{\pi}2 \) de \( \frac{\pi}4 \).}{}{}}
\End

\Data\ID{1003108204}
\Updated{2020-07-15 03:07:21}
\Level{B}
\SubArea{Definite integral}
\LANGen{
\Question{
The value of the definite integral \( \int\limits_{-\frac{\pi}4}^{\frac{\pi}4}\left(\mathrm{tg}^2x+1\right)\mathrm{d}x \) is:}
{a whole number}{a decimal number}{a proper fraction}{an irrational number}{}{}}
\LANGpl{
\Question{
Wartość całki oznaczonej \( \int\limits_{-\frac{\pi}4}^{\frac{\pi}4}\left(\mathrm{tg}^2x+1\right)\mathrm{d}x \) jest:}
{liczbą całkowitą dodatnią}{liczbą dziesiętną}{ułamkiem właściwym}{liczbą niewymierną}{}{}}
\LANGcs{
\Question{
Hodnotou určitého integrálu \( \int\limits_{-\frac{\pi}4}^{\frac{\pi}4}\left(\mathrm{tg}^2x+1\right)\mathrm{d}x \) je:}
{celé číslo}{desetinné číslo}{pravý zlomek}{iracionální číslo}{}{}}
\LANGsk{
\Question{
Hodnotou určitého integrálu \( \int\limits_{-\frac{\pi}4}^{\frac{\pi}4}\left(\mathrm{tg}^2x+1\right)\mathrm{d}x \) je:}
{celé číslo}{desatinné číslo}{pravý zlomok}{iracionálne číslo}{}{}}
\LANGes{
\Question{
El valor de la integral definida \( \int\limits_{-\frac{\pi}4}^{\frac{\pi}4}\left(\mathrm{tg}^2x+1\right)\mathrm{d}x \) es:}
{un número entero}{un número decimal}{una fracción propia}{un número irracional}{}{}}
\End

\Data\ID{1003108205}
\Updated{2020-07-15 03:07:21}
\Level{B}
\SubArea{Definite integral}
\LANGen{
\Question{
Compare the two definite integrals \( I_1=\int\limits_{-1}^1\left(x+\frac{\pi}2\right)\mathrm{d}x \) and \( I_2=\int\limits_0^{\frac{\pi}4}\mathrm{tg}\,x\cdot\cos x\,\mathrm{d}x \).}
{\( I_1 \) is bigger than \( I_2 \).}{\( I_1 \) is smaller than \( I_2 \).}{\( I_1 \) is equal to \( I_2 \).}{These integrals cannot be compared.}{}{}}
\LANGpl{
\Question{
Porównaj dwie całki oznaczone \( I_1=\int\limits_{-1}^1\left(x+\frac{\pi}2\right)\mathrm{d}x \) i \( I_2=\int\limits_0^{\frac{\pi}4}\mathrm{tg}\,x\cdot\cos x\,\mathrm{d}x \).}
{\( I_1 \) jest większa niż \( I_2 \).}{\( I_1 \) jest mniejsza niż \( I_2 \).}{\( I_1 \) jest równa \( I_2 \).}{Nie można wykonać porównania.}{}{}}
\LANGcs{
\Question{
Porovnejte dva určité integrály \( I_1=\int\limits_{-1}^1\left(x+\frac{\pi}2\right)\mathrm{d}x \) a \( I_2=\int\limits_0^{\frac{\pi}4}\mathrm{tg}\,x\cdot\cos x\,\mathrm{d}x \).}
{\( I_1 \) je větší než \( I_2 \).}{\( I_1 \) je menší než \( I_2 \).}{\( I_1 \) je roven \( I_2 \).}{Integrály nelze porovnat.}{}{}}
\LANGsk{
\Question{
Porovnajte dva určité integrály \( I_1=\int\limits_{-1}^1\left(x+\frac{\pi}2\right)\mathrm{d}x \) a \( I_2=\int\limits_0^{\frac{\pi}4}\mathrm{tg}\,x\cdot\cos x\,\mathrm{d}x \).}
{\( I_1 \) je väčší než \( I_2 \).}{\( I_1 \) je menší než \( I_2 \).}{\( I_1 \) je rovný \( I_2 \).}{Integrály sa nedajú porovnať.}{}{}}
\LANGes{
\Question{
Compara dos integrales definidas \( I_1=\int\limits_{-1}^1\left(x+\frac{\pi}2\right)\mathrm{d}x \) z \( I_2=\int\limits_0^{\frac{\pi}4}\mathrm{tg}\,x\cdot\cos x\,\mathrm{d}x \).}
{\( I_1 \) es más grande que \( I_2 \).}{\( I_1 \) es más pequeña que \( I_2 \).}{\( I_1 \) es igual a \( I_2 \).}{Estas integrales no se pueden comparar.}{}{}}
\End

\Data\ID{1003118801}
\Updated{2020-07-15 03:07:21}
\Level{B}
\SubArea{Definite integral}
\LANGen{
\Question{
Which of the following formulas is not equal to \( \int\limits_4^8\frac{3x+1}{x^2-x-6}\,\mathrm{d}x \)?}
{\( \int\limits_4^8\frac2{x-3}\,\mathrm{d}x -\int\limits_4^8\frac1{x+2}\,\mathrm{d}x \)}{\( \int\limits_4^8\frac2{x-3}\,\mathrm{d}x + \int\limits_4^8\frac1{x+2}\,\mathrm{d}x \)}{\( \int\limits_4^6\frac{3x+1}{x^2-x-6}\,\mathrm{d}x + \int\limits_6^8\frac{3x+1}{x^2-6-x}\,\mathrm{d}x \)}{\( \int\limits_8^4\frac{-1-3x}{x^2-x-6}\,\mathrm{d}x \)}{}{}}
\LANGpl{
\Question{
Który z poniższych wzorów nie jest równy \( \int\limits_4^8\frac{3x+1}{x^2-x-6}\,\mathrm{d}x \)?}
{\( \int\limits_4^8\frac2{x-3}\,\mathrm{d}x -\int\limits_4^8\frac1{x+2}\,\mathrm{d}x \)}{\( \int\limits_4^8\frac2{x-3}\,\mathrm{d}x + \int\limits_4^8\frac1{x+2}\,\mathrm{d}x \)}{\( \int\limits_4^6\frac{3x+1}{x^2-x-6}\,\mathrm{d}x + \int\limits_6^8\frac{3x+1}{x^2-6-x}\,\mathrm{d}x \)}{\( \int\limits_8^4\frac{-1-3x}{x^2-x-6}\,\mathrm{d}x \)}{}{}}
\LANGcs{
\Question{
Kterému z výrazů není roven \( \int\limits_4^8\frac{3x+1}{x^2-x-6}\,\mathrm{d}x \)?}
{\( \int\limits_4^8\frac2{x-3}\,\mathrm{d}x -\int\limits_4^8\frac1{x+2}\,\mathrm{d}x \)}{\( \int\limits_4^8\frac2{x-3}\,\mathrm{d}x + \int\limits_4^8\frac1{x+2}\,\mathrm{d}x \)}{\( \int\limits_4^6\frac{3x+1}{x^2-x-6}\,\mathrm{d}x + \int\limits_6^8\frac{3x+1}{x^2-6-x}\,\mathrm{d}x \)}{\( \int\limits_8^4\frac{-1-3x}{x^2-x-6}\,\mathrm{d}x \)}{}{}}
\LANGsk{
\Question{
Ktorému z výrazov nie je rovný \( \int\limits_4^8\frac{3x+1}{x^2-x-6}\,\mathrm{d}x \)?}
{\( \int\limits_4^8\frac2{x-3}\,\mathrm{d}x -\int\limits_4^8\frac1{x+2}\,\mathrm{d}x \)}{\( \int\limits_4^8\frac2{x-3}\,\mathrm{d}x + \int\limits_4^8\frac1{x+2}\,\mathrm{d}x \)}{\( \int\limits_4^6\frac{3x+1}{x^2-x-6}\,\mathrm{d}x + \int\limits_6^8\frac{3x+1}{x^2-6-x}\,\mathrm{d}x \)}{\( \int\limits_8^4\frac{-1-3x}{x^2-x-6}\,\mathrm{d}x \)}{}{}}
\LANGes{
\Question{
¿Cuál de las siguientes expresiones no es igual a \( \int\limits_4^8\frac{3x+1}{x^2-x-6}\,\mathrm{d}x \)?}
{\( \int\limits_4^8\frac2{x-3}\,\mathrm{d}x -\int\limits_4^8\frac1{x+2}\,\mathrm{d}x \)}{\( \int\limits_4^8\frac2{x-3}\,\mathrm{d}x + \int\limits_4^8\frac1{x+2}\,\mathrm{d}x \)}{\( \int\limits_4^6\frac{3x+1}{x^2-x-6}\,\mathrm{d}x + \int\limits_6^8\frac{3x+1}{x^2-6-x}\,\mathrm{d}x \)}{\( \int\limits_8^4\frac{-1-3x}{x^2-x-6}\,\mathrm{d}x \)}{}{}}
\End

\Data\ID{1003118802}
\Updated{2020-07-15 03:07:21}
\Level{B}
\SubArea{Definite integral}
\LANGen{
\Question{
Find  positive real numbers \( a \) and \( b \), \( a \), \( b>0 \), such that \( b-a=6 \) and \( \int\limits_a^b(4x-6)\,\mathrm{d}x=60 \).}
{\( a=1 \), \( b=7 \)}{\( a=7 \), \( b=1 \)}{\( a=24 \), \( b=30 \)}{\( a=1.5 \), \( b=7.5 \)}{}{}}
\LANGpl{
\Question{
Wskaż dodatnie liczby rzeczywiste \( a \) i \( b \), \( a \), \( b>0 \), tak, aby \( b-a=6 \) i \( \int\limits_a^b(4x-6)\,\mathrm{d}x=60 \).}
{\( a=1 \), \( b=7 \)}{\( a=7 \), \( b=1 \)}{\( a=24 \), \( b=30 \)}{\( a=1{,}5 \), \( b=7{,}5 \)}{}{}}
\LANGcs{
\Question{
Pro která kladná reálná čísla \( a \) a \( b \), taková, že \( b-a=6 \), platí \( \int\limits_a^b(4x-6)\,\mathrm{d}x=60 \)?}
{\( a=1 \), \( b=7 \)}{\( a=7 \), \( b=1 \)}{\( a=24 \), \( b=30 \)}{\( a=1{,}5 \), \( b=7{,}5 \)}{}{}}
\LANGsk{
\Question{
Pre ktoré kladné reálne čísla \( a \) a \( b \), také, že \( b-a=6 \), platí \( \int\limits_a^b(4x-6)\,\mathrm{d}x=60 \)?}
{\( a=1 \), \( b=7 \)}{\( a=7 \), \( b=1 \)}{\( a=24 \), \( b=30 \)}{\( a=1{,}5 \), \( b=7{,}5 \)}{}{}}
\LANGes{
\Question{
Halla números positivos reales \( a \) y \( b \), \( a \), \( b>0 \), tal que \( b-a=6 \) y \( \int\limits_a^b(4x-6)\,\mathrm{d}x=60 \).}
{\( a=1 \), \( b=7 \)}{\( a=7 \), \( b=1 \)}{\( a=24 \), \( b=30 \)}{\( a=1.5 \), \( b=7.5 \)}{}{}}
\End

\Data\ID{1003118804}
\Updated{2020-07-15 03:07:21}
\Level{B}
\SubArea{Definite integral}
\LANGen{
\Question{
Evaluate the integral \( \int\limits_{-3}^2 f(x)\,\mathrm{d}x \), where \( f(x)=x^{-4} \) for \( x\in\left[-3;-\frac12\right] \) and \( f(x)=12.8-6.4x \) for \( x\in\left[ -\frac12;2\right]\). (rounded to \( 2 \) decimal places)}
{\( 22.65 \)}{\( 43.\overline{8} \)}{\( 44.\overline{1} \)}{\( 29.7\overline{1} \)}{}{}}
\LANGpl{
\Question{
Wyznacz całkę \( \int\limits_{-3}^2 f(x)\,\mathrm{d}x \), gdzie \( f(x)=x^{-4} \) dla \( x\in\left\langle-3;-\frac12\right\rangle \) i \( f(x)=12{,}8-6{,}4x \) dla \( x\in\left\langle -\frac12;2\right\rangle\). (zaokrąglone do \( 2 \) miejsc po przecinku)}
{\( 22{,}65 \)}{\( 43{,}\overline{8} \)}{\( 44{,}\overline{1} \)}{\( 29{,}7\overline{1} \)}{}{}}
\LANGcs{
\Question{
Vypočítejte integrál \( \int\limits_{-3}^2 f(x)\,\mathrm{d}x \), je-li \( f(x)=x^{-4} \) pro \( x\in\left\langle-3;-\frac12\right\rangle \) a \( f(x)=12{,}8-6{,}4x \) pro \( x\in\left\langle -\frac12;2\right\rangle\). (zaokrouhleno na dvě desetinná místa)}
{\( 22{,}65 \)}{\( 43{,}\overline{8} \)}{\( 44{,}\overline{1} \)}{\( 29{,}7\overline{1} \)}{}{}}
\LANGsk{
\Question{
Vypočítajte integrál \( \int\limits_{-3}^2 f(x)\,\mathrm{d}x \), ak je \( f(x)=x^{-4} \) pre \( x\in\left\langle-3;-\frac12\right\rangle \) a \( f(x)=12{,}8-6{,}4x \) pre \( x\in\left\langle -\frac12;2\right\rangle\). (zaokrúhlené na dve desatinné miesta)}
{\( 22{,}65 \)}{\( 43{,}\overline{8} \)}{\( 44{,}\overline{1} \)}{\( 29{,}7\overline{1} \)}{}{}}
\LANGes{
\Question{
Evalúa la integral \( \int\limits_{-3}^2 f(x)\,\mathrm{d}x \), donde \( f(x)=x^{-4} \) para \( x\in\left[-3;-\frac12\right] \) y \( f(x)=12.8-6.4x \) para \( x\in\left[ -\frac12;2\right]\). (redondea a \( 2 \) decimales)}
{\( 22.65 \)}{\( 43.\overline{8} \)}{\( 44.\overline{1} \)}{\( 29.7\overline{1} \)}{}{}}
\End

\Data\ID{1003118805}
\Updated{2020-07-15 03:07:21}
\Level{B}
\SubArea{Definite integral}
\LANGen{
\Question{
Evaluate the integral \( \int\limits_{-5}^5 f(x)\,\mathrm{d}x \), where \( f(x)=x^2+2 \) for \( x\in[-5;1] \) and \( f(x)=3 \) for \( x\in[ 1;5] \).}
{\( 66 \)}{\( \frac73 \)}{\( \frac{25}3 \)}{\( 42 \)}{}{}}
\LANGpl{
\Question{
Wyznacz całkę \( \int\limits_{-5}^5 f(x)\,\mathrm{d}x \), jeśli \( f(x)=x^2+2 \) dla \( x\in\langle-5;1\rangle \) i \( f(x)=3 \) dla \( x\in\langle 1;5\rangle \).}
{\( 66 \)}{\( \frac73 \)}{\( \frac{25}3 \)}{\( 42 \)}{}{}}
\LANGcs{
\Question{
Vypočítejte integrál \( \int\limits_{-5}^5 f(x)\,\mathrm{d}x \), je-li \( f(x)=x^2+2 \) pro \( x\in\langle-5;1\rangle \) a \( f(x)=3 \) pro \( x\in\langle 1;5\rangle \).}
{\( 66 \)}{\( \frac73 \)}{\( \frac{25}3 \)}{\( 42 \)}{}{}}
\LANGsk{
\Question{
Vypočítajte integrál \( \int\limits_{-5}^5 f(x)\,\mathrm{d}x \), ak je \( f(x)=x^2+2 \) pre \( x\in\langle-5;1\rangle \) a \( f(x)=3 \) pre \( x\in\langle 1;5\rangle \).}
{\( 66 \)}{\( \frac73 \)}{\( \frac{25}3 \)}{\( 42 \)}{}{}}
\LANGes{
\Question{
Evalúa la integral \( \int\limits_{-5}^5 f(x)\,\mathrm{d}x \), donde \( f(x)=x^2+2 \) para \( x\in[-5;1] \) y \( f(x)=3 \) para \( x\in[ 1;5] \).}
{\( 66 \)}{\( \frac73 \)}{\( \frac{25}3 \)}{\( 42 \)}{}{}}
\End

\Data\ID{1003118901}
\Updated{2020-07-15 03:07:19}
\Level{B}
\SubArea{Definite integral}
\LANGen{
\Question{
Evaluate the definite integral.
\[ \int\limits_{-2}^3 \sqrt{5x+12}\,\mathrm{d}x \]}
{\( 18.3 \)}{\( 19.1 \)}{\( 18.1 \)}{\( 19.3 \)}{}{}}
\LANGpl{
\Question{
Oblicz całkę oznaczoną.
\[ \int\limits_{-2}^3 \sqrt{5x+12}\,\mathrm{d}x \]}
{\( 18{,}3 \)}{\( 19{,}1 \)}{\( 18{,}1 \)}{\( 19{,}3 \)}{}{}}
\LANGcs{
\Question{
Vypočítejte určitý integrál.
\[ \int\limits_{-2}^3 \sqrt{5x+12}\,\mathrm{d}x \]}
{\( 18{,}3 \)}{\( 19{,}1 \)}{\( 18{,}1 \)}{\( 19{,}3 \)}{}{}}
\LANGsk{
\Question{
Vypočítajte určitý integrál.
\[ \int\limits_{-2}^3 \sqrt{5x+12}\,\mathrm{d}x \]}
{\( 18{,}3 \)}{\( 19{,}1 \)}{\( 18{,}1 \)}{\( 19{,}3 \)}{}{}}
\LANGes{
\Question{
Evalúa la integral definida.
\[ \int\limits_{-2}^3 \sqrt{5x+12}\,\mathrm{d}x \]}
{\( 18.3 \)}{\( 19.1 \)}{\( 18.1 \)}{\( 19.3 \)}{}{}}
\End

\Data\ID{1003118902}
\Updated{2020-07-15 03:07:19}
\Level{B}
\SubArea{Definite integral}
\LANGen{
\Question{
Evaluate the definite integral.
\[ \int\limits_0^{10}\frac{4x}{x^2+1}\mathrm{d}x \]}
{\( \ln(10\,201) \)}{\( \ln(10\,404) \)}{\( \ln(202) \)}{\( 9 \)}{}{}}
\LANGpl{
\Question{
Oblicz całkę oznaczoną.
\[ \int\limits_0^{10}\frac{4x}{x^2+1}\mathrm{d}x \]}
{\( \ln(10\,201) \)}{\( \ln(10\,404) \)}{\( \ln(202) \)}{\( 9 \)}{}{}}
\LANGcs{
\Question{
Vypočítejte určitý integrál.
\[ \int\limits_0^{10}\frac{4x}{x^2+1}\mathrm{d}x  \]}
{\( \ln(10\,201) \)}{\( \ln(10\,404) \)}{\( \ln(202) \)}{\( 9 \)}{}{}}
\LANGsk{
\Question{
Vypočítajte určitý integrál.
\[ \int\limits_0^{10}\frac{4x}{x^2+1}\mathrm{d}x \]}
{\( \ln(10\,201) \)}{\( \ln(10\,404) \)}{\( \ln(202) \)}{\( 9 \)}{}{}}
\LANGes{
\Question{
Evalúa la integral definida.
\[ \int\limits_0^{10}\frac{4x}{x^2+1}\mathrm{d}x \]}
{\( \ln(10\,201) \)}{\( \ln(10\,404) \)}{\( \ln(202) \)}{\( 9 \)}{}{}}
\End

\Data\ID{1003118903}
\Updated{2020-07-15 03:07:19}
\Level{B}
\SubArea{Definite integral}
\LANGen{
\Question{
Evaluate the definite integral.
\[ \int\limits_{-0.5}^{0.1}\mathrm{e}^{5x-1}\mathrm{d}x \]}
{\( 0.115 \)}{\( 0.127 \)}{\( 0.121 \)}{\( 0.205 \)}{}{}}
\LANGpl{
\Question{
Oblicz całkę oznaczoną.
\[ \int\limits_{-0{,}5}^{0{,}1}\mathrm{e}^{5x-1}\mathrm{d}x \]}
{\( 0{,}115 \)}{\( 0{,}127 \)}{\( 0{,}121 \)}{\( 0{,}205 \)}{}{}}
\LANGcs{
\Question{
Vypočítejte určitý integrál.
\[ \int\limits_{-0{,}5}^{0{,}1}\mathrm{e}^{5x-1}\mathrm{d}x  \]}
{\( 0{,}115 \)}{\( 0{,}127 \)}{\( 0{,}121 \)}{\( 0{,}205 \)}{}{}}
\LANGsk{
\Question{
Vypočítajte určitý integrál.
\[ \int\limits_{-0{,}5}^{0{,}1}\mathrm{e}^{5x-1}\mathrm{d}x  \]}
{\( 0{,}115 \)}{\( 0{,}127 \)}{\( 0{,}121 \)}{\( 0{,}205 \)}{}{}}
\LANGes{
\Question{
Evalúa la integral definida.
\[ \int\limits_{-0.5}^{0.1}\mathrm{e}^{5x-1}\mathrm{d}x \]}
{\( 0.115 \)}{\( 0.127 \)}{\( 0.121 \)}{\( 0.205 \)}{}{}}
\End

\Data\ID{1003118904}
\Updated{2020-07-15 03:07:19}
\Level{B}
\SubArea{Definite integral}
\LANGen{
\Question{
Evaluate the definite integral.
\[ \int\limits_{-\frac{\pi}2}^{\pi}x\cdot\sin x\,\mathrm{d}x \]}
{\( \pi+1 \)}{\( \pi-1 \)}{\( \frac{\pi+2}2 \)}{\( \frac{\pi}2 \)}{}{}}
\LANGpl{
\Question{
Oblicz całkę oznaczoną.
\[ \int\limits_{-\frac{\pi}2}^{\pi}x\cdot\sin x\,\mathrm{d}x \]}
{\( \pi+1 \)}{\( \pi-1 \)}{\( \frac{\pi+2}2 \)}{\( \frac{\pi}2 \)}{}{}}
\LANGcs{
\Question{
Vypočítejte určitý integrál.
\[ \int\limits_{-\frac{\pi}2}^{\pi}x\cdot\sin x\,\mathrm{d}x \]}
{\( \pi+1 \)}{\( \pi-1 \)}{\( \frac{\pi+2}2 \)}{\( \frac{\pi}2 \)}{}{}}
\LANGsk{
\Question{
Vypočítajte určitý integrál.
\[ \int\limits_{-\frac{\pi}2}^{\pi}x\cdot\sin x\,\mathrm{d}x \]}
{\( \pi+1 \)}{\( \pi-1 \)}{\( \frac{\pi+2}2 \)}{\( \frac{\pi}2 \)}{}{}}
\LANGes{
\Question{
Evalúa la integral definida.
\[ \int\limits_{-\frac{\pi}2}^{\pi}x\cdot\sin x\,\mathrm{d}x \]}
{\( \pi+1 \)}{\( \pi-1 \)}{\( \frac{\pi+2}2 \)}{\( \frac{\pi}2 \)}{}{}}
\End

\Data\ID{1103118803}
\Updated{2020-07-15 03:07:21}
\Level{B}
\SubArea{Definite integral}
\IMG{1103118803.pdf}
\LANGen{
\Question{
The picture shows a graph of a power function \( f(x) \). Find the value of the definite integral of this function between \( -1 \) and \( -\frac12 \).}
{\( 1 \)}{\( 3 \)}{\( -1 \)}{\( 2 \)}{}{}}
\LANGpl{
\Question{
Na rysunku znajduje się wykres funkcji potęgowej \( f(x) \). Wyznacz wartość całki oznaczonej tej funkcji pomiędzy \( -1 \) i \( -\frac12 \).}
{\( 1 \)}{\( 3 \)}{\( -1 \)}{\( 2 \)}{}{}}
\LANGcs{
\Question{
Na obrázku je graficky zadaná mocninná funkce \( f(x) \). Vypočítejte hodnotu určitého integrálu této funkce v mezích od  \( -1 \) do \( -\frac12 \).}
{\( 1 \)}{\( 3 \)}{\( -1 \)}{\( 2 \)}{}{}}
\LANGsk{
\Question{
Na obrázku je graficky zadaná mocninná funkcia \( f(x) \). Vypočítajte hodnotu určitého integrálu tejto funkcie v intervale od  \( -1 \) do \( -\frac12 \).}
{\( 1 \)}{\( 3 \)}{\( -1 \)}{\( 2 \)}{}{}}
\LANGes{
\Question{
La imagen demuestra la gráfica de la función racional \( f(x) \). Halla el valor de la integral definida de esta función entre \( -1 \) y \( -\frac12 \).}
{\( 1 \)}{\( 3 \)}{\( -1 \)}{\( 2 \)}{}{}}
\End

\Data\ID{9000150405}
\Updated{2020-07-15 03:07:37}
\Level{B}
\SubArea{Definite integral}
\LANGen{
\Question{
Evaluate the following definite integral. 
\[ 
 \int _{-1}^{3} \frac{1}
{2x + 3}\, \text{d}x
\]}
{\(\ln 3\)}{\(\ln 1\)}{\(\ln 9\)}{\(2\ln 9\)}{}{}}
\LANGpl{
\Question{
Oblicz.
\[ 
 \int _{-1}^{3} \frac{1}
{2x + 3}\, \text{d}x
\]}
{\(\ln 3\)}{\(\ln 1\)}{\(\ln 9\)}{\(2\ln 9\)}{}{}}
\LANGcs{
\Question{
Vypočítejte \(\int _{-1}^{3} \frac{1}
{2x+3}\, \text{d}x\).}
{\(\ln 3\)}{\(\ln 1\)}{\(\ln 9\)}{\(2\ln 9\)}{}{}}
\LANGsk{
\Question{
Vypočítajte \(\int _{-1}^{3} \frac{1}
{2x+3}\, \text{d}x\).}
{\(\ln 3\)}{\(\ln 1\)}{\(\ln 9\)}{\(2\ln 9\)}{}{}}
\LANGes{
\Question{
Evalúa la siguiente integral definida.
\[ 
 \int _{-1}^{3} \frac{1}
{2x + 3}\, \text{d}x
\]}
{\(\ln 3\)}{\(\ln 1\)}{\(\ln 9\)}{\(2\ln 9\)}{}{}}
\End

\Data\ID{9000150406}
\Updated{2020-07-15 03:07:37}
\Level{B}
\SubArea{Definite integral}
\LANGen{
\Question{
Evaluate the following definite integral. 
\[ 
 \int _{-3}^{2} \frac{x}
{x^{2} + 3}\, \text{d}x
\]}
{\(\frac{1}
{2}\ln \frac{7}
{12}\)}{\(2\ln \frac{12}
 {7} \)}{\(2\ln \frac{7}
{12}\)}{\(\frac{1}
{2}\ln \frac{12}
 {7} \)}{}{}}
\LANGpl{
\Question{
Oblicz.
\[ 
 \int _{-3}^{2} \frac{x}
{x^{2} + 3}\, \text{d}x
\]}
{\(\frac{1}
{2}\ln \frac{7}
{12}\)}{\(2\ln \frac{12}
 {7} \)}{\(2\ln \frac{7}
{12}\)}{\(\frac{1}
{2}\ln \frac{12}
 {7} \)}{}{}}
\LANGcs{
\Question{
Vypočítejte \(\int _{-3}^{2} \frac{x}
{x^{2}+3}\, \text{d}x\).}
{\(\frac{1}
{2}\ln \frac{7}
{12}\)}{\(2\ln \frac{12}
 {7} \)}{\(2\ln \frac{7}
{12}\)}{\(\frac{1}
{2}\ln \frac{12}
 {7} \)}{}{}}
\LANGsk{
\Question{
Vypočítajte \(\int _{-3}^{2} \frac{x}
{x^{2}+3}\, \text{d}x\).}
{\(\frac{1}
{2}\ln \frac{7}
{12}\)}{\(2\ln \frac{12}
 {7} \)}{\(2\ln \frac{7}
{12}\)}{\(\frac{1}
{2}\ln \frac{12}
 {7} \)}{}{}}
\LANGes{
\Question{
Evalúa la siguiente integral definida.
\[ 
 \int _{-3}^{2} \frac{x}
{x^{2} + 3}\, \text{d}x
\]}
{\(\frac{1}
{2}\ln \frac{7}
{12}\)}{\(2\ln \frac{12}
 {7} \)}{\(2\ln \frac{7}
{12}\)}{\(\frac{1}
{2}\ln \frac{12}
 {7} \)}{}{}}
\End

\Data\ID{1003027501}
\Updated{2020-07-15 03:07:12}
\Level{C}
\SubArea{Definite integral}
\LANGen{
\Question{
Evaluate the following definite integral.
\[ \int_{\mathrm{e}}^{\mathrm{e}^3}\frac{10x-14}{x\cdot(x-2)}\mathrm{d}x \]}
{\( 3\ln\frac{\mathrm{e}^3-2}{\mathrm{e}-2}+14 \)}{\( 3\ln\left[\left(\mathrm{e}^3-2\right)\left(\mathrm{e}-2\right)\right]+14 \)}{\( \ln\frac{\mathrm{e}^3-2}{\mathrm{e}-2}+14 \)}{\( 3\ln\left(\mathrm{e}^3-2\right)+\ln(\mathrm{e}-2)^3+14 \)}{}{}}
\LANGpl{
\Question{
Wyznacz całkę.
\[ \int_{\mathrm{e}}^{\mathrm{e}^3}\frac{10x-14}{x\cdot(x-2)}\mathrm{d}x \]}
{\( 3\ln\frac{\mathrm{e}^3-2}{\mathrm{e}-2}+14 \)}{\( 3\ln\left[\left(\mathrm{e}^3-2\right)\left(\mathrm{e}-2\right)\right]+14 \)}{\( \ln\frac{\mathrm{e}^3-2}{\mathrm{e}-2}+14 \)}{\( 3\ln\left(\mathrm{e}^3-2\right)+\ln(\mathrm{e}-2)^3+14 \)}{}{}}
\LANGcs{
\Question{
Vypočítejte určitý integrál.
\[ \int_{\mathrm{e}}^{\mathrm{e}^3}\frac{10x-14}{x\cdot(x-2)}\mathrm{d}x \]}
{\( 3\ln\frac{\mathrm{e}^3-2}{\mathrm{e}-2}+14 \)}{\( 3\ln\left[\left(\mathrm{e}^3-2\right)\left(\mathrm{e}-2\right)\right]+14 \)}{\( \ln\frac{\mathrm{e}^3-2}{\mathrm{e}-2}+14 \)}{\( 3\ln\left(\mathrm{e}^3-2\right)+\ln(\mathrm{e}-2)^3+14 \)}{}{}}
\LANGsk{
\Question{
Vypočítajte určitý integrál.
\[ \int_{\mathrm{e}}^{\mathrm{e}^3}\frac{10x-14}{x\cdot(x-2)}\mathrm{d}x \]}
{\( 3\ln\frac{\mathrm{e}^3-2}{\mathrm{e}-2}+14 \)}{\( 3\ln\left[\left(\mathrm{e}^3-2\right)\left(\mathrm{e}-2\right)\right]+14 \)}{\( \ln\frac{\mathrm{e}^3-2}{\mathrm{e}-2}+14 \)}{\( 3\ln\left(\mathrm{e}^3-2\right)+\ln(\mathrm{e}-2)^3+14 \)}{}{}}
\LANGes{
\Question{
Evalúa la siguiente integral definida. 
\[ \int_{\mathrm{e}}^{\mathrm{e}^3}\frac{10x-14}{x\cdot(x-2)}\mathrm{d}x \]}
{\( 3\ln\frac{\mathrm{e}^3-2}{\mathrm{e}-2}+14 \)}{\( 3\ln\left[\left(\mathrm{e}^3-2\right)\left(\mathrm{e}-2\right)\right]+14 \)}{\( \ln\frac{\mathrm{e}^3-2}{\mathrm{e}-2}+14 \)}{\( 3\ln\left(\mathrm{e}^3-2\right)+\ln(\mathrm{e}-2)^3+14 \)}{}{}}
\End

\Data\ID{1003027502}
\Updated{2020-07-15 03:07:12}
\Level{C}
\SubArea{Definite integral}
\LANGen{
\Question{
Evaluate the following definite integral.
\[ \int\limits_5^{25}\frac{4x+1}{x^2-x-2}\mathrm{d}x \]}
{\( \ln\frac{13}3+3\ln\frac{23}3 \)}{\( \ln(26\cdot23^3\cdot6\cdot3^3) \)}{\(  \ln\frac{26}{27}+3\ln23+\ln6 \)}{\( \ln(24\cdot27^3)-\ln(4\cdot7^3) \)}{}{}}
\LANGpl{
\Question{
Wyznacz całkę.
\[ \int\limits_5^{25}\frac{4x+1}{x^2-x-2}\mathrm{d}x \]}
{\( \ln\frac{13}3+3\ln\frac{23}3 \)}{\( \ln(26\cdot23^3\cdot6\cdot3^3) \)}{\(  \ln\frac{26}{27}+3\ln23+\ln6 \)}{\( \ln(24\cdot27^3)-\ln(4\cdot7^3) \)}{}{}}
\LANGcs{
\Question{
Vypočítejte určitý integrál.
\[ \int\limits_5^{25}\frac{4x+1}{x^2-x-2}\mathrm{d}x \]}
{\( \ln\frac{13}3+3\ln\frac{23}3 \)}{\( \ln(26\cdot23^3\cdot6\cdot3^3) \)}{\( \ln\frac{26}{27}+3\ln23+\ln6 \)}{\( \ln(24\cdot27^3)-\ln(4\cdot7^3) \)}{}{}}
\LANGsk{
\Question{
Vypočítajte určitý integrál.
\[ \int\limits_5^{25}\frac{4x+1}{x^2-x-2}\mathrm{d}x \]}
{\( \ln\frac{13}3+3\ln\frac{23}3 \)}{\( \ln(26\cdot23^3\cdot6\cdot3^3) \)}{\(  \ln\frac{26}{27}+3\ln23+\ln6 \)}{\( \ln(24\cdot27^3)-\ln(4\cdot7^3) \)}{}{}}
\LANGes{
\Question{
Evalúa la siguiente integral definida. 
\[ \int\limits_5^{25}\frac{4x+1}{x^2-x-2}\mathrm{d}x \]}
{\( \ln\frac{13}3+3\ln\frac{23}3 \)}{\( \ln(26\cdot23^3\cdot6\cdot3^3) \)}{\(  \ln\frac{26}{27}+3\ln23+\ln6 \)}{\( \ln(24\cdot27^3)-\ln(4\cdot7^3) \)}{}{}}
\End

\Data\ID{1003027503}
\Updated{2020-07-15 03:07:12}
\Level{C}
\SubArea{Definite integral}
\LANGen{
\Question{
Evaluate the following definite integral.
\[ \int\limits_{\mathrm{e}}^{12}\frac{x+5}{\frac14 x\cdot(x+4)}\mathrm{d}x \]}
{\(\ln(\mathrm{e}+4)+5\ln12-2\ln4-5 \)}{\( 5+2\ln 4+5\ln12-\ln(\mathrm{e}+4) \)}{\( \ln\frac{15}4-5+\ln(4+\mathrm{e}) \)}{\( \frac{ \ln \frac{12^5}{16} }{ \ln \frac{\mathrm{e}^5}{\mathrm{e}+4}} \)}{}{}}
\LANGpl{
\Question{
Wyznacz całkę.
\[ \int\limits_{\mathrm{e}}^{12}\frac{x+5}{\frac14 x\cdot(x+4)}\mathrm{d}x \]}
{\( \ln(\mathrm{e}+4)+5\ln12-2\ln4-5 \)}{\( 5+2\ln 4+5\ln12-\ln(\mathrm{e}+4) \)}{\(\ln\frac{15}4-5+\ln(4+\mathrm{e}) \)}{\( \frac{ \ln \frac{12^5}{16} }{ \ln \frac{\mathrm{e}^5}{\mathrm{e}+4}} \)}{}{}}
\LANGcs{
\Question{
Vypočítejte určitý integrál.
\[ \int\limits_{\mathrm{e}}^{12}\frac{x+5}{\frac14 x\cdot(x+4)}\mathrm{d}x \]}
{\( \ln(\mathrm{e}+4)+5\ln12-2\ln4-5 \)}{\( 5+2\ln 4+5\ln12-\ln(\mathrm{e}+4) \)}{\( \ln\frac{15}4-5+\ln(4+\mathrm{e}) \)}{\( \frac{ \ln \frac{12^5}{16} }{ \ln \frac{\mathrm{e}^5}{\mathrm{e}+4}} \)}{}{}}
\LANGsk{
\Question{
Vypočítajte určitý integrál.
\[ \int\limits_{\mathrm{e}}^{12}\frac{x+5}{\frac14 x\cdot(x+4)}\mathrm{d}x \]}
{\( \ln(\mathrm{e}+4)+5\ln12-2\ln4-5 \)}{\( 5+2\ln 4+5\ln12-\ln(\mathrm{e}+4) \)}{\( \ln\frac{15}4-5+\ln(4+\mathrm{e}) \)}{\( \frac{ \ln \frac{12^5}{16} }{ \ln \frac{\mathrm{e}^5}{\mathrm{e}+4}} \)}{}{}}
\LANGes{
\Question{
Evalúa la siguiente integral definida. 
\[ \int\limits_{\mathrm{e}}^{12}\frac{x+5}{\frac14 x\cdot(x+4)}\mathrm{d}x \]}
{\(\ln(\mathrm{e}+4)+5\ln12-2\ln4-5 \)}{\( 5+2\ln 4+5\ln12-\ln(\mathrm{e}+4) \)}{\( \ln\frac{15}4-5+\ln(4+\mathrm{e}) \)}{\( \frac{ \ln \frac{12^5}{16} }{ \ln \frac{\mathrm{e}^5}{\mathrm{e}+4}} \)}{}{}}
\End

\Data\ID{1003118905}
\Updated{2020-07-15 03:07:19}
\Level{C}
\SubArea{Definite integral}
\LANGen{
\Question{
Evaluate the definite integral.
\[ \int\limits_{\mathrm{e}}^{\mathrm{e}^2}\frac{\ln x}x\,\mathrm{d}x \]}
{\( 1.5 \)}{\( 2.5 \)}{\( \mathrm{e}^2 \)}{\( \mathrm{e}^2-\mathrm{e} \)}{}{}}
\LANGpl{
\Question{
Oblicz całkę oznaczoną.
\[ \int\limits_{\mathrm{e}}^{\mathrm{e}^2}\frac{\ln x}x\,\mathrm{d}x \]}
{\( 1{,}5 \)}{\( 2{,}5 \)}{\( \mathrm{e}^2 \)}{\( \mathrm{e}^2-\mathrm{e} \)}{}{}}
\LANGcs{
\Question{
Vypočítejte určitý integrál.
\[ \int\limits_{\mathrm{e}}^{\mathrm{e}^2}\frac{\ln x}x\,\mathrm{d}x \]}
{\( 1{,}5 \)}{\( 2{,}5 \)}{\( \mathrm{e}^2 \)}{\( \mathrm{e}^2-\mathrm{e} \)}{}{}}
\LANGsk{
\Question{
Vypočítajte určitý integrál.
\[ \int\limits_{\mathrm{e}}^{\mathrm{e}^2}\frac{\ln x}x\,\mathrm{d}x \]}
{\( 1{,}5 \)}{\( 2{,}5 \)}{\( \mathrm{e}^2 \)}{\( \mathrm{e}^2-\mathrm{e} \)}{}{}}
\LANGes{
\Question{
Evalúa la integral definida.
\[ \int\limits_{\mathrm{e}}^{\mathrm{e}^2}\frac{\ln x}x\,\mathrm{d}x \]}
{\( 1.5 \)}{\( 2.5 \)}{\( \mathrm{e}^2 \)}{\( \mathrm{e}^2-\mathrm{e} \)}{}{}}
\End

\Data\ID{1003118906}
\Updated{2020-07-15 03:07:19}
\Level{C}
\SubArea{Definite integral}
\LANGen{
\Question{
Evaluate the definite integral.
\[ \int\limits_{1}^{\mathrm{e}}\left(\sqrt[3]x-x^3\right)\cdot\ln x\,\mathrm{d}x \]}
{\( -9.03 \)}{\( 9.03 \)}{\( 11.4 \)}{\( -11.4 \)}{}{}}
\LANGpl{
\Question{
Oblicz całkę oznaczoną.
\[ \int\limits_{1}^{\mathrm{e}}\left(\sqrt[3]x-x^3\right)\cdot\ln x\,\mathrm{d}x \]}
{\( -9{,}03 \)}{\( 9{,}03 \)}{\( 11{,}4 \)}{\( -11{,}4 \)}{}{}}
\LANGcs{
\Question{
Vypočítejte určitý integrál.
\[ \int\limits_{1}^{\mathrm{e}}\left(\sqrt[3]x-x^3\right)\cdot\ln x\,\mathrm{d}x \]}
{\( -9{,}03 \)}{\( 9{,}03 \)}{\( 11{,}4 \)}{\( -11{,}4 \)}{}{}}
\LANGsk{
\Question{
Vypočítajte určitý integrál.
\[ \int\limits_{1}^{\mathrm{e}}\left(\sqrt[3]x-x^3\right)\cdot\ln x\,\mathrm{d}x \]}
{\( -9{,}03 \)}{\( 9{,}03 \)}{\( 11{,}4 \)}{\( -11{,}4 \)}{}{}}
\LANGes{
\Question{
Evalúa la integral definida.
\[ \int\limits_{1}^{\mathrm{e}}\left(\sqrt[3]x-x^3\right)\cdot\ln x\,\mathrm{d}x \]}
{\( -9.03 \)}{\( 9.03 \)}{\( 11.4 \)}{\( -11.4 \)}{}{}}
\End

\Data\ID{1003124302}
\Updated{2020-07-15 03:07:19}
\Level{C}
\SubArea{Definite integral}
\LANGen{
\Question{
Which of the given values of a real number \( a\in\left(\frac{\pi}2;\pi\right) \) makes the equality \( \int\limits_a^{2a} (3\sin x-4x)\,\mathrm{d}x=-6a^2 \) true?}
{\( a =\frac23\pi \)}{\( a=\frac56\pi \)}{\( a=\frac34\pi \)}{\( a=\frac78\pi \)}{}{}}
\LANGpl{
\Question{
Która z podanych wartości liczby rzeczywistej \( a\in\left(\frac{\pi}2;\pi\right) \) sprawia, że równość  \( \int\limits_a^{2a} (3\sin x-4x)\,\mathrm{d}x=-6a^2 \) jest prawdziwa?}
{\( a =\frac23\pi \)}{\( a=\frac56\pi \)}{\( a=\frac34\pi \)}{\( a=\frac78\pi \)}{}{}}
\LANGcs{
\Question{
Pro které reálné číslo \( a\in\left(\frac{\pi}2;\pi\right) \) platí rovnost \( \int\limits_a^{2a} (3\sin x-4x)\,\mathrm{d}x=-6a^2 \)?}
{\( a =\frac23\pi \)}{\( a=\frac56\pi \)}{\( a=\frac34\pi \)}{\( a=\frac78\pi \)}{}{}}
\LANGsk{
\Question{
Pre ktoré reálne číslo \( a\in\left(\frac{\pi}2;\pi\right) \) platí rovnosť \( \int\limits_a^{2a} (3\sin x-4x)\,\mathrm{d}x=-6a^2 \)?}
{\( a =\frac23\pi \)}{\( a=\frac56\pi \)}{\( a=\frac34\pi \)}{\( a=\frac78\pi \)}{}{}}
\LANGes{
\Question{
¿Para qué número real \( a\in\left(\frac{\pi}2;\pi\right) \) vale la igualdad \( \int\limits_a^{2a} (3\sin x-4x)\,\mathrm{d}x=-6a^2 \)?}
{\( a =\frac23\pi \)}{\( a=\frac56\pi \)}{\( a=\frac34\pi \)}{\( a=\frac78\pi \)}{}{}}
\End

\Data\ID{1003124303}
\Updated{2020-07-15 03:07:19}
\Level{C}
\SubArea{Definite integral}
\LANGen{
\Question{
Which of the given values of real numbers \( a \), \( b\in\left(0;\frac{\pi}2\right) \), such as \( a < b \), makes the equality \( \int\limits_a^b \cos x\,\mathrm{d}x=2\cos\frac{\pi}4\cdot\sin\frac{\pi}{12} \) true?}
{\( a=\frac{\pi}6 \), \( b=\frac{\pi}3 \)}{\( a=\frac{\pi}3 \), \( b=\frac{\pi}6 \)}{\( a=\frac{\pi}3 \), \( b=\frac{\pi}4 \)}{\( a=\frac{\pi}4 \), \( b=\frac{\pi}3 \)}{}{}}
\LANGpl{
\Question{
Która z podanych wartości liczby rzeczywistej \( a \), \( b\in\left(0;\frac{\pi}2\right) \) tak, aby \( a < b \), sprawia, że równość  \( \int\limits_a^b \cos x\,\mathrm{d}x=2\cos\frac{\pi}4\cdot\sin\frac{\pi}{12} \) jest prawdziwa?}
{\( a=\frac{\pi}6 \), \( b=\frac{\pi}3 \)}{\( a=\frac{\pi}3 \), \( b=\frac{\pi}6 \)}{\( a=\frac{\pi}3 \), \( b=\frac{\pi}4 \)}{\( a=\frac{\pi}4 \), \( b=\frac{\pi}3 \)}{}{}}
\LANGcs{
\Question{
Pro která reálná čísla \( a \), \( b\in\left(0;\frac{\pi}2\right) \) taková, že \( a < b \), platí rovnost \( \int\limits_a^b \cos x\,\mathrm{d}x=2\cos\frac{\pi}4\cdot\sin\frac{\pi}{12} \)?}
{\( a=\frac{\pi}6 \), \( b=\frac{\pi}3 \)}{\( a=\frac{\pi}3 \), \( b=\frac{\pi}6 \)}{\( a=\frac{\pi}3 \), \( b=\frac{\pi}4 \)}{\( a=\frac{\pi}4 \), \( b=\frac{\pi}3 \)}{}{}}
\LANGsk{
\Question{
Pre ktorá reálne čísla \( a \), \( b\in\left(0;\frac{\pi}2\right) \) také, že \( a < b \), platí rovnosť \( \int\limits_a^b \cos x\,\mathrm{d}x=2\cos\frac{\pi}4\cdot\sin\frac{\pi}{12} \)?}
{\( a=\frac{\pi}6 \), \( b=\frac{\pi}3 \)}{\( a=\frac{\pi}3 \), \( b=\frac{\pi}6 \)}{\( a=\frac{\pi}3 \), \( b=\frac{\pi}4 \)}{\( a=\frac{\pi}4 \), \( b=\frac{\pi}3 \)}{}{}}
\LANGes{
\Question{
¿Para qué número real \( a \), \( b\in\left(0;\frac{\pi}2\right) \), tal que \( a < b \), vale la igualdad \( \int\limits_a^b \cos x\,\mathrm{d}x=2\cos\frac{\pi}4\cdot\sin\frac{\pi}{12} \)?}
{\( a=\frac{\pi}6 \), \( b=\frac{\pi}3 \)}{\( a=\frac{\pi}3 \), \( b=\frac{\pi}6 \)}{\( a=\frac{\pi}3 \), \( b=\frac{\pi}4 \)}{\( a=\frac{\pi}4 \), \( b=\frac{\pi}3 \)}{}{}}
\End

\Data\ID{1003124304}
\Updated{2020-07-15 03:07:12}
\Level{C}
\SubArea{Definite integral}
\LANGen{
\Question{
Given a function \( f(x)=ax^4+bx \), find real numbers \( a \) and \( b \), such that \( \int\limits_0^1f(x)\,\mathrm{d}x=27 \) and \( \int\limits_{-1}^0f(x)\,\mathrm{d}x=57 \).}
{\( a=210 \), \( b=-30 \)}{\( a=210 \), \( b=30 \)}{\( a=75 \), \( b=60 \)}{\( a=30 \), \( b=210 \)}{}{}}
\LANGpl{
\Question{
Dana jest funkcja \( f(x)=ax^4+bx \), wskaż liczby rzeczywiste \( a \) i \( b \) tak, aby \( \int\limits_0^1f(x)\,\mathrm{d}x=27 \) i \( \int\limits_{-1}^0f(x)\,\mathrm{d}x=57 \).}
{\( a=210 \), \( b=-30 \)}{\( a=210 \), \( b=30 \)}{\( a=75 \), \( b=60 \)}{\( a=30 \), \( b=210 \)}{}{}}
\LANGcs{
\Question{
Pro funkci \( f(x)=ax^4+bx \) najděte taková reálná čísla \( a \) a \( b \), aby platilo \( \int\limits_0^1f(x)\,\mathrm{d}x=27 \) a \( \int\limits_{-1}^0f(x)\,\mathrm{d}x=57 \).}
{\( a=210 \), \( b=-30 \)}{\( a=210 \), \( b=30 \)}{\( a=75 \), \( b=60 \)}{\( a=30 \), \( b=210 \)}{}{}}
\LANGsk{
\Question{
Pre funkciu \( f(x)=ax^4+bx \) nájdite také reálna čísla \( a \) a \( b \), aby platilo \( \int\limits_0^1f(x)\,\mathrm{d}x=27 \) a \( \int\limits_{-1}^0f(x)\,\mathrm{d}x=57 \).}
{\( a=210 \), \( b=-30 \)}{\( a=210 \), \( b=30 \)}{\( a=75 \), \( b=60 \)}{\( a=30 \), \( b=210 \)}{}{}}
\LANGes{
\Question{
Dada la función \( f(x)=ax^4+bx \), halla números reales \( a \) y \( b \), tal que \( \int\limits_0^1f(x)\,\mathrm{d}x=27 \) y \( \int\limits_{-1}^0f(x)\,\mathrm{d}x=57 \).}
{\( a=210 \), \( b=-30 \)}{\( a=210 \), \( b=30 \)}{\( a=75 \), \( b=60 \)}{\( a=30 \), \( b=210 \)}{}{}}
\End

\Data\ID{1003124305}
\Updated{2020-07-15 03:07:19}
\Level{C}
\SubArea{Definite integral}
\LANGen{
\Question{
Given a function \( f(x)=ax^6+bx^3+cx+8 \), find real numbers \( a \), \( b \) and \( c \), such that \( \int\limits_0^1f(x)\,\mathrm{d}x=\frac{35}4 \), \( f'(0)=2 \) and \( f'(1)=180 \).}
{\( a=7 \), \( b=-5 \), \( c=2 \)}{\( a=7 \), \( b=5 \), \( c=2 \)}{\( a=-7 \), \( b=-5 \), \( c=2 \)}{\( a=-7 \), \( b=5 \), \( c=-2 \)}{}{}}
\LANGpl{
\Question{
Dana jest funkcja \( f(x)=ax^6+bx^3+cx+8 \), wskaż liczby rzeczywiste  \( a \), \( b \) i  \( c \) tak, aby  \( \int\limits_0^1f(x)\,\mathrm{d}x=\frac{35}4 \), \( f'(0)=2 \) i \( f'(1)=180 \).}
{\( a=7 \), \( b=-5 \), \( c=2 \)}{\( a=7 \), \( b=5 \), \( c=2 \)}{\( a=-7 \), \( b=-5 \), \( c=2 \)}{\( a=-7 \), \( b=5 \), \( c=-2 \)}{}{}}
\LANGcs{
\Question{
Pro funkci \( f(x)=ax^6+bx^3+cx+8 \) najděte taková reálná čísla \( a \), \( b \) a \( c \), aby platilo \( \int\limits_0^1f(x)\,\mathrm{d}x=\frac{35}4 \), \( f'(0)=2 \) a \( f'(1)=180 \).}
{\( a=7 \), \( b=-5 \), \( c=2 \)}{\( a=7 \), \( b=5 \), \( c=2 \)}{\( a=-7 \), \( b=-5 \), \( c=2 \)}{\( a=-7 \), \( b=5 \), \( c=-2 \)}{}{}}
\LANGsk{
\Question{
Pre funkciu \( f(x)=ax^6+bx^3+cx+8 \) nájdite také reálna čísla \( a \), \( b \) a \( c \), aby platilo \( \int\limits_0^1f(x)\,\mathrm{d}x=\frac{35}4 \), \( f'(0)=2 \) a \( f'(1)=180 \).}
{\( a=7 \), \( b=-5 \), \( c=2 \)}{\( a=7 \), \( b=5 \), \( c=2 \)}{\( a=-7 \), \( b=-5 \), \( c=2 \)}{\( a=-7 \), \( b=5 \), \( c=-2 \)}{}{}}
\LANGes{
\Question{
Dada la función \( f(x)=ax^6+bx^3+cx+8 \), halla números reales \( a \), \( b \) y \( c \), tal que \( \int\limits_0^1f(x)\,\mathrm{d}x=\frac{35}4 \), \( f'(0)=2 \) y \( f'(1)=180 \).}
{\( a=7 \), \( b=-5 \), \( c=2 \)}{\( a=7 \), \( b=5 \), \( c=2 \)}{\( a=-7 \), \( b=-5 \), \( c=2 \)}{\( a=-7 \), \( b=5 \), \( c=-2 \)}{}{}}
\End

\Data\ID{1003124306}
\Updated{2020-07-15 03:07:19}
\Level{C}
\SubArea{Definite integral}
\LANGen{
\Question{
Which of the positive values of a real number \(a\) makes the equality \( \int\limits_1^a \frac3{3x+5}\,\mathrm{d}x=\ln\frac74 \) true?}
{\( a = 3 \)}{\( a = 2 \)}{\( a = 4 \)}{\( a = 5 \)}{}{}}
\LANGpl{
\Question{
Która z dodatnich  wartości liczby rzeczywistej \(a\) sprawia, że równość  \( \int\limits_1^a \frac3{3x+5}\,\mathrm{d}x=\ln\frac74 \) jest prawdziwa?}
{\( a = 3 \)}{\( a = 2 \)}{\( a = 4 \)}{\( a = 5 \)}{}{}}
\LANGcs{
\Question{
Pro které kladné reálné číslo \(a\) platí rovnost \( \int\limits_1^a \frac3{3x+5}\,\mathrm{d}x=\ln\frac74 \)?}
{\( a = 3 \)}{\( a = 2 \)}{\( a = 4 \)}{\( a = 5 \)}{}{}}
\LANGsk{
\Question{
Pre ktoré kladné reálne číslo \(a\) platí rovnosť \( \int\limits_1^a \frac3{3x+5}\,\mathrm{d}x=\ln\frac74 \)?}
{\( a = 3 \)}{\( a = 2 \)}{\( a = 4 \)}{\( a = 5 \)}{}{}}
\LANGes{
\Question{
¿Para qué número real positivo \(a\) vale la igualdad \( \int\limits_1^a \frac3{3x+5}\,\mathrm{d}x=\ln\frac74 \)?}
{\( a = 3 \)}{\( a = 2 \)}{\( a = 4 \)}{\( a = 5 \)}{}{}}
\End

\Data\ID{1003124307}
\Updated{2020-07-15 03:07:19}
\Level{C}
\SubArea{Definite integral}
\LANGen{
\Question{
Which of the positive values of a real number \(a\) makes the equality \(\int\limits_a^7 \frac{6x-5}{3x^2-5x}\,\mathrm{d}x=\ln 56 \) true?}
{\( a=2 \)}{\( a=1 \)}{\( a=3 \)}{\( a=5 \)}{}{}}
\LANGpl{
\Question{
Która z dodatnich wartości liczby rzeczywistej \(a\) sprawia, że równość  \(\int\limits_a^7 \frac{6x-5}{3x^2-5x}\,\mathrm{d}x=\ln 56 \) jest prawdziwa?}
{\( a=2 \)}{\( a=1 \)}{\( a=3 \)}{\( a=5 \)}{}{}}
\LANGcs{
\Question{
Pro které kladné reálné číslo \(a\) platí rovnost \(\int\limits_a^7 \frac{6x-5}{3x^2-5x}\,\mathrm{d}x=\ln 56 \)?}
{\( a=2 \)}{\( a=1 \)}{\( a=3 \)}{\( a=5 \)}{}{}}
\LANGsk{
\Question{
Pre ktoré kladné reálne číslo \(a\) platí rovnosť \(\int\limits_a^7 \frac{6x-5}{3x^2-5x}\,\mathrm{d}x=\ln 56 \)?}
{\( a=2 \)}{\( a=1 \)}{\( a=3 \)}{\( a=5 \)}{}{}}
\LANGes{
\Question{
¿Para qué número real positivo \(a\) vale la igualdad \(\int\limits_a^7 \frac{6x-5}{3x^2-5x}\,\mathrm{d}x=\ln 56 \)?}
{\( a=2 \)}{\( a=1 \)}{\( a=3 \)}{\( a=5 \)}{}{}}
\End

\Data\ID{1003124308}
\Updated{2020-07-15 03:07:12}
\Level{C}
\SubArea{Definite integral}
\LANGen{
\Question{
Which of the values of a real number \( a\in(\pi;2\pi) \) makes the equality \( \int\limits_{\pi}^{a+\pi} \sin2x\,\mathrm{d}x=1 \) true?}
{\( a=\frac32\pi \)}{\( a=\frac34\pi \)}{\( a=\frac12\pi \)}{\( a=\frac23\pi \)}{}{}}
\LANGpl{
\Question{
Która z wartości liczby rzeczywistej \( a\in(\pi;2\pi) \) sprawia, że równość  \( \int\limits_{\pi}^{a+\pi} \sin2x\,\mathrm{d}x=1 \) jest prawdziwa?}
{\( a=\frac32\pi \)}{\( a=\frac34\pi \)}{\( a=\frac12\pi \)}{\( a=\frac23\pi \)}{}{}}
\LANGcs{
\Question{
Pro které reálné číslo \( a\in(\pi;2\pi) \) platí rovnost \( \int\limits_{\pi}^{a+\pi} \sin2x\,\mathrm{d}x=1 \)?}
{\( a=\frac32\pi \)}{\( a=\frac34\pi \)}{\( a=\frac12\pi \)}{\( a=\frac23\pi \)}{}{}}
\LANGsk{
\Question{
Pre ktoré reálne číslo \( a\in(\pi;2\pi) \) platí rovnosť \( \int\limits_{\pi}^{a+\pi} \sin2x\,\mathrm{d}x=1 \)?}
{\( a=\frac32\pi \)}{\( a=\frac34\pi \)}{\( a=\frac12\pi \)}{\( a=\frac23\pi \)}{}{}}
\LANGes{
\Question{
¿Para qué número real \( a\in(\pi;2\pi) \) vale la igualdad \( \int\limits_{\pi}^{a+\pi} \sin2x\,\mathrm{d}x=1 \)?}
{\( a=\frac32\pi \)}{\( a=\frac34\pi \)}{\( a=\frac12\pi \)}{\( a=\frac23\pi \)}{}{}}
\End

\Data\ID{1103124301}
\Updated{2020-07-15 03:07:19}
\Level{C}
\SubArea{Definite integral}
\IMG{1103124301.pdf}
\LANGen{
\Question{
The picture shows graphs of two quadratic functions \( f_1(x) \) and \( f_2(x) \). Find the unknown real positive constant \( a \) (as shown in the picture) such that the value of the definite integral \( \int\limits_{-1}^1 f_1(x)\,\mathrm{d}x \) is greater by \( 8 \) than the value of the definite integral \( \int\limits_{-1}^1 f_2(x)\,\mathrm{d}x \).}
{\( a = 3 \)}{\( a = 1 \)}{\( a = 4 \)}{\( a = 6 \)}{}{}}
\LANGpl{
\Question{
Rysunek przedstawia wykresy dwóch funkcji kwadratowych  \( f_1(x) \) i \( f_2(x) \). Wskaż wartość stałej  \( a \) (spójrz na rysunek) tak, aby wartość całki oznaczonej l \( \int\limits_{-1}^1 f_1(x)\,\mathrm{d}x \) była większa o \( 8 \) od wartości całki oznaczonej  \( \int\limits_{-1}^1 f_2(x)\,\mathrm{d}x \).}
{\( a = 3 \)}{\( a = 1 \)}{\( a = 4 \)}{\( a = 6 \)}{}{}}
\LANGcs{
\Question{
Na obrázku jsou grafy dvou kvadratických funkcí  \( f_1(x) \) a \( f_2(x) \). Určete neznámou kladnou reálnou  konstantu \( a \) z obrázku tak, aby byla hodnota určitého integrálu \( \int\limits_{-1}^1 f_1(x)\,\mathrm{d}x \) o \( 8 \) větší než hodnota určitého integrálu \( \int\limits_{-1}^1 f_2(x)\,\mathrm{d}x \).}
{\( a = 3 \)}{\( a = 1 \)}{\( a = 4 \)}{\( a = 6 \)}{}{}}
\LANGsk{
\Question{
Na obrázku sú grafy dvoch kvadratických funkcií  \( f_1(x) \) a \( f_2(x) \). Určte neznámu kladnú reálnu  konštantu \( a \) z obrázku tak, aby bola hodnota určitého integrálu \( \int\limits_{-1}^1 f_1(x)\,\mathrm{d}x \) o \( 8 \) väčšia než hodnota určitého integrálu \( \int\limits_{-1}^1 f_2(x)\,\mathrm{d}x \).}
{\( a = 3 \)}{\( a = 1 \)}{\( a = 4 \)}{\( a = 6 \)}{}{}}
\LANGes{
\Question{
La imagen muestra gráficas de dos funciones cuadráticas \( f_1(x) \) y \( f_2(x) \). Halla la constante real positiva desconocida \( a \) (tal como está en la imagen) tal que el valor de la integral definida \( \int\limits_{-1}^1 f_1(x)\,\mathrm{d}x \) es más grande de \( 8 \) que el valor de la integral definida \( \int\limits_{-1}^1 f_2(x)\,\mathrm{d}x \).}
{\( a = 3 \)}{\( a = 1 \)}{\( a = 4 \)}{\( a = 6 \)}{}{}}
\End
